%% Article : Comparative genomics of animal vs human MTBC lineages
%% Structure initiale — 2026-03-26

\documentclass[11pt,a4paper]{article}

\usepackage[utf8]{inputenc}
\usepackage[T1]{fontenc}
\usepackage[english]{babel}
\usepackage{amsmath, amssymb}
\usepackage{graphicx}
\usepackage{booktabs}
\usepackage{longtable}
\usepackage{geometry}
\usepackage{hyperref}
\usepackage{authblk}
\usepackage{xcolor}
\usepackage{multirow}
\usepackage{siunitx}
\usepackage{caption}
\usepackage{subcaption}
\usepackage[numbers,sort&compress]{natbib}

\geometry{margin=2.5cm}

% Helpers
\newcommand{\Mtb}{\textit{M.~tuberculosis}}
\newcommand{\MTBCfull}{\textit{Mycobacterium tuberculosis} complex (MTBC)}
\newcommand{\MTBC}{MTBC}
\newcommand{\Msur}{\textit{M.~suricattae}}
\newcommand{\Mmun}{\textit{M.~mungi}}
\newcommand{\Mdas}{\textit{M.~dassievlei}}
\newcommand{\Mchi}{\textit{M.~chimpanzae}}
\newcommand{\todo}[1]{\textbf{\textcolor{red}{[TODO: #1]}}}

% ─────────────────────────────────────────────────────────────────────────────
\title{Phylogenomic signatures of multi-host adaptation in the \MTBC:\\
Convergent PKS/PDIM remodelling in animal lineages\\
and secondary human re-adaptation of L6/L9/L10}

\author[1]{Christophe Guyeux}
\affil[1]{FEMTO-ST Institute, Univ.\ Bourgogne Franche-Comté, CNRS, Supmicrotech,
          Belfort, France}

\date{\today}
% ─────────────────────────────────────────────────────────────────────────────

\begin{document}

\maketitle

% ─────────────────────────────────────────────────────────────────────────────
\begin{abstract}
The \MTBCfull{} encompasses both human-adapted lineages responsible for the global
tuberculosis pandemic and a set of animal-adapted lineages occupying restricted
ecological niches.
%
We present the first large-scale comparative genomic analysis contrasting the four
basal animal-adapted lineages --- \Msur{} ($n = 17$), \textit{M.~dassievlei} ($n = 7$),
\Mmun{} ($n = 1$), and \Mchi{} ($n = 1$) --- against secondarily human-adapted lineages
(L6, L9, L10; $n = 1{,}132$ combined) and strict human-adapted lineages
(L1--L4; $n \geq 5$ representatives each).
%
From a dataset of \num{274800} MTBC genomes accessed through TBannotator~v3.6,
we characterised 904 SPDIs (sequence variant identifiers) differentiating the animal
clade core genome from its secondarily human-adapted sister clade (L6/L9/L10).
Ancestry stratification reveals that 68\% of these are ancestral MTBC variants
(H37Rv reference bias); only 23~SPDIs (2.5\%) are truly exclusive to
the four animal lineages, including \textit{eccC3}~D1146N (ESX-3) and
\textit{rpoA}~V325F (RNA polymerase~$\alpha$).
%
Analysis reveals a 25-fold higher insertion sequence (IS) activity in animal
versus human lineages (2.1\% vs 0.0\%), with IS6110 preferentially targeting
\textit{PPE55} (68.4\% of total IS activity). Functional annotation revealed a significant enrichment of PE/PPE family mutations
in animal lineages ($p = 0.0002$ vs.\ ancestral core), alongside a catalogue of
40 animal-differentiating loss-of-function (LoF) mutations.
%
All five ESX type-VII secretion systems harbour animal-differentiating missense variants,
predominantly in ESX-5 ($n = 9$ SPDIs), yet none exhibit LoF, indicating
functional remodelling rather than pathway loss.
%
The PKS/PDIM lipid biosynthesis pathway shows the strongest signal of adaptive
divergence: 98 SPDIs across 24 genes, with missense/synonymous ratios
(dN/dS proxy) exceeding 2.0 in \textit{ppsA}, \textit{pks8}, and
\textit{pks13}, and convergent pseudogenisation of
\textit{pks15} in all four animal lineages.
%
Convergence analysis reveals 333 genes independently mutated in all four animal
species (8.3\% of the genome; confirmed in the two multi-genome lineages, $n \geq 7$), representing the
largest convergent gene set documented in bacterial pathogen evolution and prominently
enriched for polyketide synthase genes.
%
Phylogenetic clade signature analysis identifies 10~synapomorphic SPDIs
defining the grand clade (L5~+ animal lineages~+ L6/L9/L10), including
functionally impactful frameshifts in \textit{ald} (L-alanine dehydrogenase)
and \textit{cyp142} (cholesterol-metabolising cytochrome~P450), alongside
14~L6/L9/L10-exclusive SPDIs marking secondary human adaptation.
Notably, no synapomorphic variants could be identified for the strict
human clade (L1--L4--L7), reflecting the reference bias inherent to SNP
calling against the L4-derived H37Rv genome.
%
These results delineate the genomic boundaries between animal and human MTBC
ecotypes and highlight the PKS/PDIM pathway as a key determinant of host range.

\noindent\textbf{Keywords:} \textit{Mycobacterium tuberculosis} complex, animal-adapted
lineages, comparative genomics, PDIM, PE/PPE, ESX, convergent evolution, host adaptation.
\end{abstract}

% ─────────────────────────────────────────────────────────────────────────────
\section{Introduction}
\label{sec:intro}

The \MTBC{} represents a remarkably homogeneous species complex, sharing
$>$99.9\% nucleotide identity at the whole-genome level, yet encompassing
pathogens with dramatically different host ranges.
%
Human-adapted lineages (L1--L4 and L7) account for the vast majority of the
estimated 10 million annual tuberculosis cases worldwide \cite{WHO2023}, while a
set of phylogenetically basal lineages --- collectively termed the ``animal
clade'' --- cause tuberculosis-like disease in specific non-human mammalian hosts
and are rarely, if ever, transmitted to humans \cite{Boritsch2016, Gagneux2025}.
%
Understanding the genomic basis of this host-range dichotomy is a central
question in MTBC evolutionary biology, with implications for the origins of the
human epidemic and for the surveillance of zoonotic risk. Key systems implicated
in host adaptation include the PE/PPE gene family, which undergoes frequent
homologous recombination and contributes to antigenic variability
\cite{Karboul2008}, the ESX type-VII secretion systems, which have evolved
sophisticated mechanisms to counter host immunity \cite{Groschel2016}, and
insertion sequence (IS) elements such as IS6110, which drive genomic plasticity
and have been linked to strain virulence and epidemiological success
\cite{Fallow2010}.

The animal-adapted MTBC encompasses four named lineages: \Msur{}
(meerkats, \textit{Suricata suricatta}) \cite{Alexander2010},
\Mmun{} (banded mongooses, \textit{Mungos mungo}) \cite{Roos2012},
\Mdas{} (rock hyraxes, \textit{Procavia capensis}) \cite{Parsons2008}, and
\Mchi{} (chimpanzees, \textit{Pan troglodytes}) \cite{Fabre2010}.
%
Recent phylogenetic analyses reveal that MTBC evolution involves a complex
pattern of multiple host-jumping events rather than simple animal-human
bifurcation. The lineages L6 (\textit{M.~africanum} West Africa~II),
L9 (\textit{M.~africanum} East Africa), and L10 \cite{Coscolla2021, Guyeux2024L10}
represent secondarily human-adapted lineages that evolved from an
animal-adapted intermediate, forming a paraphyletic relationship with
the primarily animal-adapted clade and predominantly restricted to
sub-Saharan Africa.
%
This complex phylogenetic topology places these lineages in a pivotal role
for understanding multi-directional host-switching and the genomic signatures
of secondary host adaptation \cite{Comas2013, Boritsch2016}.

Lineage 5 (L5, \textit{M.~africanum} West Africa~I) occupies a transitional
phylogenetic position that represents the primary human-adapted branch
diverging before the animal adaptation event that later gave rise to both
the primarily animal-adapted lineages and the secondarily human-adapted
L6/L9/L10 clade \cite{deJong2010, Coll2014}.
%
This complex evolutionary scenario, supported by phylogenetic evidence from
thousands of genome-scale analyses, provides an informative framework for
reconstructing the sequence of multiple host-jumping events that shaped
modern MTBC diversity through recurrent adaptation and counter-adaptation
across mammalian host species.

Previous genomic studies of the animal-adapted MTBC have focused principally on
\textit{M.~bovis} and \textit{M.~caprae}, which were excluded from the present
analysis as they belong to a distinct phylogenetic subgroup \cite{Brosch2002}.
%
For the four primarily animal-adapted lineages considered here,
characterisation has been limited to the definition of species-specific
Regions of Difference (RD) \cite{Parsons2008, Roos2012, Fabre2010},
spoligotype profiles, and broad phylogenetic placement \cite{Coll2014, Boritsch2016, Brites2018}.
%
No study has systematically compared the single-nucleotide variant landscape of
all four primarily animal-adapted lineages against both the strict human-adapted
lineages (L1--L4) and the secondarily human-adapted lineages (L6/L9/L10) at
scale, nor attempted a quantitative analysis of functional pathway divergence,
loss-of-function catalogues, or convergent evolution patterns that distinguish
primary versus secondary host adaptation events.

Here, leveraging a dataset of \num{274800} MTBC genomes catalogued in
TBannotator~v3.6 \cite{Guyeux2024}, we perform the first large-scale comparative
genomic analysis between primarily animal-adapted lineages and both strict
human-adapted lineages (L1--L4) and secondarily human-adapted lineages (L6/L9/L10).
%
Our study addresses five main questions:
(i)~What is the shared core genomic identity of the primarily animal-adapted clade?
(ii)~Which functional systems are differentially remodelled between primary
animal adaptation and secondary human adaptation?
(iii)~Is there evidence for convergent evolution across independently evolving
animal lineages?
(iv)~What molecular signatures distinguish secondary human adaptation (L6/L9/L10)
from primary human adaptation (L1--L4)?
(v)~What constitutes the minimal universal genomic signature shared by all
human-adapted MTBC lineages regardless of their evolutionary origin?

% ─────────────────────────────────────────────────────────────────────────────
\section{Materials and Methods}
\label{sec:methods}

\subsection{Genome dataset and lineage definitions}
\label{sec:methods:data}

Whole-genome sequences and associated metadata were accessed through TBannotator~v3.6,
a curated database of \num{274800} \textit{M.~tuberculosis} complex genomes with
standardised SPDI-format variant calls referenced to H37Rv (NC\_000962.3)
\cite{Guyeux2024}.

Five groups were defined for comparative analysis
(Table~\ref{tab:groups}):
%
\begin{enumerate}
  \item \textbf{Animal clade} (\textit{Suricattae} $n = 17$;
        \textit{Mungi} $n = 1$; \textit{Dassie} $n = 7$;
        \textit{Chimpanzé} $n = 1$; total $n = 26$).
  \item \textbf{Secondarily human-adapted clade} (L6 $n = 1{,}089$; L9 $n = 40$;
        L10 $n = 3$; total $n = 1{,}132$).
  \item \textbf{Strict human lineages} (L1--L4 and L7; five representative
        strains per sub-lineage for L1--L4, H37Rv excluded;
        L7 $n = 81$).
  \item \textbf{L5} ($n = 443$; basal position within the grand clade).
  \item \textbf{Ancestral core}: intersection of animal and secondarily human-adapted
        clades ($n = 388$ SPDIs).
\end{enumerate}

\subsection{Core SPDI analysis}
\label{sec:methods:spdi}

For each group, the core genome was defined as the set of SPDIs present
in 100\% of member strains (intersection).
%
Animal-differentiating SPDIs were defined as core animal SPDIs absent from
the secondarily human-adapted core (L6/L9/L10 intersection).
%
To assess the true exclusivity of these variants, each animal-differentiating
SPDI was cross-referenced against all available MTBC lineages: L1--L4
(including $>$50~L4 sub-lineages; 5~representative strains each,
excluding H37Rv), L5 ($n = 5$), L7 ($n = 5$), and \textit{M.~canettii}
($n = 10$). SPDIs present in \textit{M.~canettii} were classified as
ancestral; those absent from \textit{canettii} but present in any strict
human lineage were classified as shared derived; the remainder were
classified as truly animal-exclusive.
%
Conversely, human-exclusive SPDIs were defined as secondarily human-adapted core SPDIs
absent from all animal lineages.
%
Phylogenetic clade signatures were computed by identifying SPDIs shared
by all members of a given clade and absent from all outgroup lineages
(tolerance: $\leq 1$ false positive per 5~sampled strains to account
for noise and potential misclassification).
Four signatures were tested: (A)~strict human clade (L1--L7--L2--L3--L4,
excluding L5); (B)~extended human clade (L1--L7--L2--L3--L4~+ L5);
(C)~grand clade (L5~+ animal lineages~+ L6/L9/L10, absent from
L1--L4, L7 and \textit{M.~canettii});
(D)~L5-specific (absent from all other lineages).

\subsection{Functional annotation}
\label{sec:methods:annotation}

Variant effects were predicted with SnpEff \cite{Cingolani2012} against
H37Rv (NC\_000962.3).
%
Functional categories were assigned according to the Mycobrowser
annotation schema \cite{Kapopoulou2011}.
%
Loss-of-function mutations (LoF) were defined as frameshifts, nonsense
(stop-gain), and start-lost variants.
%
A dN/dS proxy was computed per gene as the raw ratio of missense to synonymous
SPDI counts, without normalisation by the number of synonymous and
non-synonymous sites available in each codon. This metric does not
constitute a true $\omega$ (dN/dS) estimate and is used as a comparative
ordinal indicator only; values derived from small counts
(e.g.\ 2~missense~/~1~synonymous $= 2.0$) carry wide implicit confidence
intervals and should not be over-interpreted.

\subsection{ESX secretion system classification}
\label{sec:methods:esx}

ESX system membership was assigned to each SPDI based on gene name and locus
tag, using a curated dictionary mapping all component genes of ESX-1 through
ESX-5 to their respective system \cite{Houben2012}.

\subsection{PKS/PDIM pathway analysis}
\label{sec:methods:pks}

Genes of the phthiocerol dimycocerosate (PDIM) biosynthesis pathway were
organised into six functional stages: precursor activation (\textit{fadD28}),
phthiocerol synthesis (\textit{ppsA--E}), mycocerosate synthesis (\textit{mas}),
ABC transport (\textit{drrA/B/C}), membrane transport (\textit{mmpL7}),
and surface translocation (\textit{lppX}), plus 14 auxiliary PKS genes
\cite{Trivedi2005, Kolly2017}.

\subsection{Convergent evolution analysis}
\label{sec:methods:convergence}

For each species pair, species-specific SPDIs were defined by subtracting the
animal clade core (904 SPDIs fixed in all four species) from each
species' individual core.
%
Convergent genes were defined as genes harbouring independent (non-identical)
species-specific SPDIs in two or more animal lineages.
%
Statistical enrichment of functional categories among convergent genes was
assessed with Fisher's exact test (two-tailed; $\alpha = 0.05$).
%
To test whether the observed convergence exceeds neutral expectation, a
permutation-based null model was constructed: for each species, its
species-specific SPDI count was redistributed randomly across genes
with probability proportional to gene length (from H37Rv GFF3 annotation),
and the number of genes hit in all four species was tallied
(10{,}000 replicates; significance assessed by $Z$-score and empirical
$p$-value).

\subsection{Phylogenetic reconstruction}
\label{sec:methods:phylo}

A neighbour-joining (NJ) tree was constructed from pairwise Jaccard distances
computed on binary SPDI profiles.
%
The dataset comprised all 26 animal strains and a subsample of 19 secondarily human-adapted
strains (5--7 per lineage).
%
Tree visualisation was performed with ETE3 \cite{Huerta-Cepas2016}.

\subsection{Robustness analysis}
\label{sec:methods:robust}

Given that \Mmun{} and \Mchi{} are each represented by a single genome,
all primary analyses were repeated restricting the animal group to the
two multi-genome lineages (\Msur{}, $n = 17$; \Mdas{}, $n = 7$).
%
Concordance with the full four-species results is reported for each major finding.

% ─────────────────────────────────────────────────────────────────────────────
\section{Results}
\label{sec:results}

\subsection{Dataset overview and phylogenetic relationships}
\label{sec:results:phylo}

The NJ tree based on SPDI Jaccard distances recapitulates the expected
MTBC topology (Figure~\ref{fig:phylo}).
%
Within the animal clade, \Msur{} and \Mdas{} form a tight cluster, while
\Mmun{} is more diverged and \Mchi{} occupies an outgroup position relative
to the three African animal lineages.
%
The secondarily human-adapted clade shows L10 as the deepest branch, followed by L9, then L6.
%
L5 is confirmed as intermediate between the animal/secondarily human-adapted grade and the
strict human-adapted L1--L4 lineages.

\subsection{Core genome signatures of the animal clade}
\label{sec:results:core}

The animal clade core comprises \num{1292} SPDIs, of which 904 are absent
from the secondarily human-adapted core (L6/L9/L10 intersection;
Table~\ref{tab:core_summary}).
%
The ancestral core shared between animal and secondarily human-adapted clades totals 388~SPDIs.
%
Conversely, only 14 SPDIs are exclusive to the secondarily human-adapted clade, nine of which
are absent from any animal lineage.

\label{sec:results:stratification}
\paragraph{Ancestry stratification of the 904 animal-differentiating SPDIs.}
Cross-referencing the 904 animal core SPDIs against all available MTBC lineages
(L1--L4 with $>$50 sub-lineages, L5, L7, and \textit{M.~canettii}; 5~representative
strains per sub-lineage) reveals that the majority are ancestral MTBC variants
shared with other lineages rather than animal-specific innovations
(Table~\ref{tab:animal_stratification}).
%
Of the 904, \num{616} (68\%) are present in \textit{M.~canettii} and represent
ancient polymorphisms at positions where the H37Rv reference carries a derived
allele---the same reference-bias phenomenon identified for the clade
signatures (Section~\ref{sec:results:clade_sig}).
%
An additional 42 are post-\textit{canettii} variants shared with strict human
lineages (L1--L4, L7).
%
After removing all variants found in any non-animal MTBC lineage, only
\textbf{23~SPDIs} (2.5\%) are truly exclusive to the animal clade
(Table~\ref{tab:truly_exclusive}).
%
Among these 23 animal-specific synapomorphies, notable variants include
\textit{eccC3}~p.D1146N (ESX-3 ATPase), \textit{rpoA}~p.V325F
(RNA polymerase $\alpha$ subunit), \textit{Rv1816}~p.G195D (transcriptional
regulator), \textit{hsaB}~p.A168G (cholesterol degradation), and
\textit{Rv2297}~p.C118* (the only loss-of-function among the 23).
%
Although the functional and enrichment analyses presented below are computed
on the full 904-SPDI set (which validly contrasts the animal core against
its phylogenetic sister clade, L6/L9/L10), these results should be
interpreted as characterising the animal core genome composition rather
than animal-specific adaptations \emph{per se}.

\begin{table}[h]
\centering
\caption{Ancestry stratification of the 904 animal core SPDIs (absent from
the L6/L9/L10 core). Each SPDI was tested against 5~representative strains
per sub-lineage across all MTBC lineages.}
\label{tab:animal_stratification}
\begin{tabular}{@{}lrr@{}}
\toprule
\textbf{Category} & \textbf{SPDIs} & \textbf{\%} \\
\midrule
Ancestral (present in \textit{M.~canettii}) & 616 & 68.1 \\
Derived, shared with strict human (L1--L4, L7) & 42 & 4.6 \\
Derived, absent strict human, shared with L5 & 54 & 6.0 \\
Derived, absent strict human \& L5, in some L6/L9/L10 & 169 & 18.7 \\
\textbf{Truly animal-exclusive} & \textbf{23} & \textbf{2.5} \\
\bottomrule
\end{tabular}
\end{table}

\begin{table}[h]
\centering
\caption{The 23 truly animal-exclusive SPDIs: present in all four animal
lineages, absent from all tested human lineages (L1--L4, L5, L7),
L6/L9/L10, and \textit{M.~canettii}.}
\label{tab:truly_exclusive}
\begin{tabular}{@{}llll@{}}
\toprule
\textbf{Gene} & \textbf{Locus} & \textbf{Effect} & \textbf{Change} \\
\midrule
\textit{eccC3}  & Rv0284  & missense  & D1146N \\
\textit{rpoA}   & Rv3457c & missense  & V325F \\
\textit{Rv1816} & Rv1816  & missense  & G195D \\
\textit{hsaB}   & Rv3567c & missense  & A168G \\
\textit{eno}    & Rv1023  & missense  & G293V \\
\textit{Rv2509} & Rv2509  & missense  & S7G \\
\textit{Rv2554c}& Rv2554c & missense  & Q165K \\
\textit{Rv2670c}& Rv2670c & missense  & P124L \\
\textit{TB22.2} & Rv3036c & missense  & V217M \\
\textit{Rv0480c}& Rv0480c & missense  & G64R \\
\textit{Rv0556} & Rv0556  & missense  & A135V \\
\textit{Rv0613c}& Rv0613c & missense  & S309F \\
\textit{Rv2297} & Rv2297  & nonsense  & C118* \\
\textit{mmpL11} & Rv0202c & synonymous & T806T \\
\textit{nrdB}   & Rv0233  & synonymous & D200D \\
\textit{pepR}   & Rv2782c & synonymous & V198V \\
\textit{Rv1597} & Rv1597  & synonymous & G101G \\
\textit{Rv1026} & Rv1026  & synonymous & P270P \\
\textit{Rv2969c}& Rv2969c & synonymous & P9P \\
\textit{PPE60}  & Rv3478  & synonymous & S302S \\
---             & ---     & intergenic & --- \\
---             & ---     & intergenic & --- \\
---             & ---     & intergenic & --- \\
\bottomrule
\end{tabular}
\end{table}

Functional categorisation of the 904 animal-differentiating SPDIs reveals conserved hypothetical proteins as the largest category (265 variants, 29\%), followed by intermediary metabolism \& respiration (146, 16\%), cell wall \& cell processes (141, 16\%), and intergenic regions (148, 16\%). PE/PPE family proteins, despite representing only 7\% of animal-differentiating variants (60 SPDIs), show significant enrichment compared to human lineages. The most heavily mutated individual loci include \textit{PPE55} (13 SPDIs), \textit{ppsA} (5 SPDIs), and \textit{PPE54} (5 SPDIs), alongside mycobacterial membrane proteins \textit{mmpL2} and metabolic genes \textit{gyrA} (4 SPDIs each). This distribution suggests that animal adaptation predominantly involves fine-tuning of core metabolic functions and host-pathogen interaction systems rather than wholesale pathway disruption.

\subsection{Functional annotation: PE/PPE enrichment and loss-of-function}
\label{sec:results:functional}

Among the 904 animal-differentiating SPDIs, functional annotation reveals a significant
enrichment of PE/PPE family mutations compared to both the secondarily human-adapted exclusive
set and the ancestral core ($p = 0.02$ and $p = 0.0002$, respectively; Fisher's
exact test; Figure~\ref{fig:enrichment}).
%
PPE55 alone carries 13 animal-differentiating SPDIs.

The animal clade harbours 40 LoF mutations in its differentiating core,
compared to only 5 in the ancestral core (Figure~\ref{fig:lof}).
%
\begin{table}[h]
\centering
\caption{Top loss-of-function mutations in the animal core genome}
\label{tab:lof_genes}
\begin{tabular}{@{}lllp{6cm}@{}}
\toprule
\textbf{Gene} & \textbf{Mutation} & \textbf{Category} & \textbf{Predicted functional consequence} \\
\midrule
\textit{mmpL12} & S1108* & Cell wall processes & Loss of specialized lipid transport; reduced cell envelope integrity \\
\textit{fadD11} & frameshift & Lipid metabolism & Impaired fatty acid activation; altered lipid biosynthesis \\
\textit{ctpI} & frameshift & Cell wall processes & Disrupted cation homeostasis; altered metal tolerance \\
\textit{PPE55} & frameshift & PE/PPE & Loss of host-pathogen interaction; reduced antigenic variability \\
\textit{PPE58} & frameshift & PE/PPE & Impaired immune evasion; altered surface antigen presentation \\
\textit{aceAa} & frameshift & Central metabolism & Reduced glyoxylate cycle activity; metabolic adaptation \\
\textit{cobL} & frameshift & Cofactor biosynthesis & Vitamin B12 synthesis defect; altered metabolic flux \\
\textit{cobK} & frameshift & Cofactor biosynthesis & Vitamin B12 synthesis defect; altered metabolic flux \\
\bottomrule
\end{tabular}
\end{table}

The 40 animal-differentiating LoF mutations predominantly affect conserved hypothetical proteins (16 mutations), followed by intermediary metabolism genes (8 mutations) and PE/PPE family proteins (6 mutations). Notably, several LoF mutations target metabolic adaptation systems (vitamin B12 biosynthesis, fatty acid metabolism) and host-interaction mechanisms (PE/PPE proteins, membrane transporters), suggesting that gene loss contributes to ecological specialization in animal-adapted lineages.

\subsection{ESX type-VII secretion systems}
\label{sec:results:esx}

All five ESX systems carry animal-differentiating SPDIs (Table~\ref{tab:esx}).
%
ESX-5, the largest and most recently acquired secretion system with key roles
in PE/PPE protein translocation \cite{Houben2012}, is the most affected, with
nine animal-differentiating SPDIs across \textit{esxK}, \textit{esxL}, \textit{esxO},
\textit{eccB5}, \textit{eccC5}, and \textit{eccD5}.
%
Critically, none of the 24 animal-differentiating ESX SPDIs are LoF: all are missense
or synonymous variants, indicating functional remodelling rather than pathway loss.
%
In contrast, human lineages carry two ESX-1 SPDIs and three ESX-2 SPDIs,
suggesting a distinct trajectory of ESX system divergence.

\subsection{PKS/PDIM pathway: selective remodelling of lipid biosynthesis}
\label{sec:results:pks}

The PDIM biosynthesis pathway is the most extensively diverged metabolic
system between animal and human ecotypes (Figure~\ref{fig:pks_overview}).
%
Across 24 PKS/PDIM genes, we identified 98 unique SPDIs: 20 in the
animal ancestral core and 78 species-specific.

The central PDIM pathway (ppsA--E, mas, fadD28, mmpL7, lppX, drrABC)
shows zero LoF mutations, with 12 missense variants concentrated in \textit{ppsA}
(dN/dS proxy $= 2.0$) and \textit{ppsE} ($= 0.5$), suggesting functional
remodelling of phthiocerol chain composition rather than pathway inactivation
(Figure~\ref{fig:pks_pathway}).
%
The ABC transporter DrrA carries a fixed H309D missense shared by all four
animal lineages.

Auxiliary PKS genes present a strikingly different pattern.
%
\textit{pks15} is pseudogenised by a frameshift (G:GC) in the core of all four
animal species, consistent with the known \textit{pks15/1} fusion distinction
between \Mtb{} and \textit{M.~bovis} \cite{Reed2004}.
%
Beyond this shared event, independent LoF mutations have arisen in
\textit{pks5} (frameshift in \Msur{} and \Mmun{}),
\textit{pks6} (frameshift in \Msur{}), and
\textit{pks7} (nonsense p.Q1265* in \Mchi{}).
%
Elevated dN/dS proxies in \textit{pks8} (2.0) and \textit{pks13} (2.0)
suggest ongoing positive or relaxed selection in these auxiliary enzymes.
%
Crucially, no PKS or PDIM gene carries any SPDI exclusive to human lineages,
confirming that this pathway is under strict purifying selection in
human-adapted MTBC (Figure~\ref{fig:pks_overview}D).

\subsection{KEGG pathway enrichment}
\label{sec:results:kegg}

KEGG pathway enrichment analysis reveals distinct metabolic signatures between animal and human MTBC lineages. Animal-differentiating genes show significant representation in oxidative phosphorylation (9 genes, 4.39\%), lipoarabinomannan biosynthesis (5 genes, 2.44\%), and cutin/suberine/wax biosynthesis (3 genes, 1.46\%), suggesting enhanced cell wall modification capacity. Conversely, secondarily human-adapted lineages are enriched for amino acid biosynthesis pathways (15 vs 0 genes in animal-differentiating; $p = 0.032$), particularly cysteine and methionine metabolism, indicating greater metabolic versatility in amino acid synthesis. The ancestral core genome shows strong enrichment for one-carbon folate metabolism (7 genes vs 2 in animal-differentiating; $p = 0.006$), DNA replication machinery, and sulfur relay systems, representing conserved essential functions maintained across both ecotypes.

Pathway completeness analysis demonstrates that core metabolic functions remain largely intact across both ecotypes (Table~\ref{tab:pathway_completeness}). Essential pathways including TCA cycle, fatty acid metabolism, mycolic acid biosynthesis, and oxidative phosphorylation show 100\% gene integrity in both animal and human lineages. However, animal lineages exhibit selective loss-of-function in peripheral pathways: glycolysis/gluconeogenesis (96.8\% intact; \textit{aldehyde dehydrogenase} Rv0223c frameshift) and cofactor biosynthesis (99.2\% intact; vitamin B12 synthesis genes \textit{cobK} and \textit{cobL} frameshifts). This metabolic rewiring pattern supports the hypothesis that animal adaptation involves specialization toward host-specific niches with targeted loss of metabolic redundancy while preserving essential virulence-associated pathways.

\begin{table}[h]
\centering
\caption{Metabolic pathway completeness in animal vs human MTBC lineages.
Total gene counts per pathway are from the KEGG database (release 2025;
organism code \texttt{mtu}, reference strain H37Rv).
LoF = loss-of-function variants (frameshift or nonsense) in the animal or
human core SPDI sets}
\label{tab:pathway_completeness}
\begin{tabular}{@{}lrrrrr@{}}
\toprule
\textbf{Pathway} & \textbf{Total} & \multicolumn{2}{c}{\textbf{Animal lineages}} & \multicolumn{2}{c}{\textbf{Human lineages}} \\
 & \textbf{genes} & \textbf{LoF} & \textbf{\% intact} & \textbf{LoF} & \textbf{\% intact} \\
\midrule
TCA cycle & 29 & 0 & 100.0 & 0 & 100.0 \\
Glycolysis/Gluconeogenesis & 31 & 1 & 96.8 & 0 & 100.0 \\
Fatty acid metabolism & 35 & 0 & 100.0 & 0 & 100.0 \\
Mycolic acid biosynthesis & 36 & 0 & 100.0 & 0 & 100.0 \\
Oxidative phosphorylation & 49 & 0 & 100.0 & 0 & 100.0 \\
One carbon pool by folate & 19 & 0 & 100.0 & 0 & 100.0 \\
Glyoxylate metabolism & 39 & 0 & 100.0 & 0 & 100.0 \\
Biosynthesis of amino acids & 111 & 0 & 100.0 & 0 & 100.0 \\
Purine metabolism & 51 & 0 & 100.0 & 0 & 100.0 \\
Pyrimidine metabolism & 27 & 0 & 100.0 & 0 & 100.0 \\
Biosynthesis of cofactors & 133 & 2 & 99.2 & 0 & 100.0 \\
Sulfur metabolism & 22 & 0 & 100.0 & 0 & 100.0 \\
\bottomrule
\end{tabular}
\end{table}

\subsection{Convergent evolution across animal lineages}
\label{sec:results:convergence}

After removing the shared animal core (904 SPDIs fixed in all four species),
333 genes carry independent (non-identical) mutations in all four animal
lineages, and 118 additional genes are convergently mutated in three of the
four species (Figure~\ref{fig:convergence}). This represents 8.3\% of the
estimated 4{,}000-gene MTBC genome.
%
To assess whether this level of convergence could arise by chance, we performed
a permutation test distributing species-specific mutations randomly across
genes proportionally to gene length (10{,}000 replicates).
The expected number of 4-of-4 convergent genes under the null model is
$27.7 \pm 4.7$, yielding a 12-fold enrichment over expectation
($Z = 65.4$; $p < 10^{-100}$; null range 13--48).
This far exceeds levels reported in other bacterial pathogens (typically
50--200 genes; 1.3--5.0\% of the genome \cite{Marvig2015}).
%
These convergently evolving genes are significantly enriched for lipid
metabolism functions ($p < 0.001$; Fisher's exact test), particularly the
polyketide synthase superfamily, with 13 hypermutated loci ($\geq$10 SPDIs each):
\textit{pks12} (17 SPDIs), \textit{PPE54} (17), \textit{PPE5} (15), \textit{pks2} (13),
\textit{ppsE} (12), and \textit{pks6} (12). Additionally, 49 genes show
convergent loss-of-function mutations across lineages, indicating diversifying
selection pressure targeting specific functional systems.
%
Beyond the dominant PKS/lipid metabolism signal, convergent evolution extends to key regulatory and metabolic systems. Signal transduction pathways show remarkable convergence, with \textit{pstP} (serine/threonine protein phosphatase; 9 independent SPDIs) and \textit{tcrY} (two-component response regulator; 9 SPDIs) independently mutated across all four animal lineages. These proteins regulate stress response and virulence networks, suggesting shared selective pressures toward modified host-interaction strategies. PE/PPE family proteins also emerge as convergent targets, with \textit{PPE54} (17 SPDIs), \textit{PPE5} (15 SPDIs), and \textit{PPE55} (11 SPDIs) showing extensive independent diversification. Metabolic convergence includes \textit{fadH} (NADPH-dependent 2,4-dienoyl-CoA reductase; 9 SPDIs), linking to fatty acid degradation, and transcriptional regulators (Rv1453, 9 SPDIs), indicating coordinated rewiring of regulatory networks. This pattern demonstrates that animal adaptation involves systematic remodelling of host-sensing, metabolic regulation, and antigenic variation systems beyond the well-characterized lipid biosynthesis pathway alterations.
%
The robustness of the convergence signal was verified by restricting the analysis
to the two multi-genome lineages (\Msur{} and \Mdas{}): 10 genes remain
convergently mutated with independent SPDIs, all within the PKS/lipid category.

\subsection{Differential IS element activity drives genomic instability}
\label{sec:results:is_elements}

Analysis of insertion sequence (IS) element activity reveals a striking 25-fold
difference between animal and human MTBC lineages (Figure~\ref{fig:is_activity}).
Animal-differentiating mutations include 19 IS-related variants (2.1\% of total),
compared to zero IS activity in secondarily human-adapted lineages (0.0\% across 223 total variants).
%
The animal clade shows preferential activity of IS6110 (13 mutations), IS1538
(4 mutations), and IS3-family elements (1 mutation), with two loss-of-function
consequences.

A pronounced IS insertion hotspot emerges at \textit{PPE55} (Rv3347c), which
harbours all 13 IS6110-associated mutations identified in the animal clade
across a 6{,}755~bp genomic span.
%
These mutations include one frameshift disruption (generating a truncated protein),
six missense variants (L2259P, F1003D, N992T, M999I, L998F, S922G), and six
synonymous changes, suggesting both ongoing transposition events and selection
pressure specifically targeting this PE/PPE locus.
%
The \textit{PPE55} hotspot represents 68.4\% of total animal IS activity and
21.7\% of all PE/PPE mutations in animal lineages, indicating that IS6110
preferentially disrupts antigenic variation systems in animal-adapted lineages.
%
Among 30 PE/PPE genes affected by animal-differentiating mutations, \textit{PPE55}
shows 2.6-fold more mutations than the next most targeted locus (\textit{PPE54}
with 5 mutations), demonstrating exceptional vulnerability to IS6110 insertion,
consistent with earlier reports of multiple independent IS6110 insertions at
identical positions in PPE genes \cite{McEvoy2009}.
%
A secondary hotspot at the IS1538 transposase gene (Rv3428c) shows evidence of
self-targeting, with four mutations including one frameshift that potentially
limits further transposition activity.

This differential IS burden suggests that animal-adapted MTBC lineages experience
ongoing genomic instability driven by mobile genetic elements, contrasting with
the genomic stability observed in human-adapted lineages. The preferential targeting
of PE/PPE genes may contribute to the antigenic diversification previously observed
in these lineages, providing a mechanistic basis for the convergent evolution of
host-pathogen interaction systems.

\subsection{Drug resistance gene validation}
\label{sec:results:dr}

Cross-referencing the animal-differentiating core against the WHO 2022 catalogue of
mutations associated with drug resistance identified 29~unique SPDIs falling
within canonical DR loci (\textit{katG}, \textit{rpoB}, \textit{embB},
\textit{embC}, \textit{embA}, \textit{gyrA}, \textit{gyrB}, \textit{kasA},
\textit{tlyA}, \textit{gid}, \textit{fgd1}, \textit{pncA}, \textit{mmpL5}).
%
None of these variants is classified as associated with resistance (WHO
Grades~1--2). Fifteen are synonymous (Grade~5), seven are classified as
``not associated with resistance'' (Grade~4), and five are recognised
lineage markers (Grade~5). Two missense variants (\textit{embC}~P772Q,
\textit{kasA}~G258S) are absent from the WHO catalogue entirely. The only
variant of uncertain significance is \textit{pncA}~P115S (Grade~3), found
exclusively in the single \Mchi{} genome.
%
Notably, the canonical animal-clade polymorphisms \textit{katG}~R463L and
\textit{embB}~E378A, shared by all four animal lineages, are confirmed as
lineage markers unrelated to isoniazid or ethambutol resistance (Grade~4).
%
The \textit{gyrA}~S95T substitution falls within the quinolone
resistance-determining region (QRDR) but is a well-characterised
animal-clade marker with no fluoroquinolone resistance association.
%
Overall, these results confirm that the extensive genomic divergence of
animal-adapted MTBC does not affect canonical drug targets in a manner
predictive of resistance.

\subsection{Comparative genomics of the two animal clades}
\label{sec:results:inter_clade}

The MTBC animal lineages comprise two phylogenetically distinct clades:
clade~1 (\textit{M.~bovis}, \textit{M.~caprae}, \textit{M.~orygis},
La4, \textit{M.~microti}, \textit{M.~pinnipedii};
$n = 71 + 178 + 218 + 21 + 29 + 12 = 529$) and clade~2 (\Msur{}, \Mdas{},
\Mmun{}, \Mchi{}; $n = 26$), which diverged independently within the
grand clade.
%
To determine what unites and distinguishes these two clades, we computed
core SPDIs for each (intersection of all available strains per species)
and compared them.

The core genomes of clade~1 (Bovis $\cap$ Caprae $\cap$ Orygis $\cap$
La4 $\cap$ Microti $\cap$ Pinnipedii; 427~SPDIs) and clade~2 (Suricattae $\cap$ Dassie
$\cap$ Mungi $\cap$ Chimpanzé; \num{1292}~SPDIs) share 413~SPDIs (97\% of
clade~1, 32\% of clade~2).
%
However, this shared set does not represent an animal-specific identity:
virtually all shared SPDIs are present in at least one L6/L9/L10 strain.
%
Strikingly, \textbf{zero SPDIs are shared by both animal clades while
being absent from all non-animal lineages}.
%
The shared core between the two animal clades consists entirely of the
grand-clade ancestral background.

In contrast, each clade harbours its own derived synapomorphies.
%
After filtering against all non-animal lineages (\textit{M.~canettii},
L1--L4, L5, L7, L6/L9/L10), clade~1 retains \textbf{2~truly exclusive
SPDIs} and clade~2 retains \textbf{23~truly exclusive SPDIs}
(Table~\ref{tab:truly_exclusive}).
%
No gene is targeted by exclusive variants in both clades simultaneously
--- their synapomorphies affect entirely non-overlapping gene sets
(Table~\ref{tab:inter_clade}).

\begin{table}[h]
\centering
\caption{Comparison of truly exclusive synapomorphies between the two
animal clades. Notable genes listed for each clade; no convergent gene
targeting is observed.}
\label{tab:inter_clade}
\begin{tabular}{@{}p{3.5cm}p{5.5cm}p{5.5cm}@{}}
\toprule
 & \textbf{Clade~1} (Bovis/Caprae/Orygis/La4/Microti/Pinnipedii) & \textbf{Clade~2} (Suricattae/Dassie/Mungi/Chimpanzé) \\
\midrule
Truly exclusive SPDIs & 2 & 23 \\
LoF & 0 & 1 (\textit{Rv2297}~C118*) \\
\addlinespace
TA systems & \textit{vapB15}~N63D (antitoxin) & --- \\
ESX secretion & --- & \textit{eccC3}~D1146N \\
Transcription & --- & \textit{rpoA}~V325F, \textit{Rv1816}~G195D \\
Cholesterol catabolism & --- & \textit{hsaB}~A168G \\
Other & 1~intergenic & \textit{eno}~G293V, \textit{Rv2509}~S7G, \textit{Rv2554c}~Q165K, + others \\
Convergent genes & \multicolumn{2}{c}{\textbf{0} --- entirely non-overlapping targets} \\
\bottomrule
\end{tabular}
\end{table}

\subsection{Relationship between clade~2 and the secondarily human-adapted lineages}
\label{sec:results:c2_vs_l6l9l10}

Clade~2 (Suricattae/Dassie/Mungi/Chimpanzé) and L6/L9/L10 are phylogenetic
sisters within the grand clade.
%
Their core genomes share 388~SPDIs, of which 255~(66\%) are ancestral
(\textit{M.~canettii}-positive), 121~(31\%) are derived variants also present
in clade~1 (grand-clade level), and 12~(3\%) are \textbf{synapomorphies
exclusive to the clade~2\,+\,L6/L9/L10 node} --- absent from clade~1, L5,
all strict human lineages, and \textit{M.~canettii}
(Table~\ref{tab:synapo_c2_l6l9l10}).

All 12 node synapomorphies show perfect concordance: 5/5~strains positive in
each of L6.1, L9, L10, Suricattae, and Dassie, and 1/1 in Mungi and
Chimpanzé, with 0/5 in every clade~1 species and L5.
%
Among the 12, seven are missense variants and two target pathways repeatedly
implicated in MTBC adaptation: \textit{eccB3}~(Rv0283; I312M), a component of
the ESX-3 type~VII secretion system involved in iron acquisition, and
\textit{plsB2}~(Rv2482c; D16G), a glycerol-3-phosphate acyltransferase in
membrane lipid biosynthesis.
%
The gene \textit{Rv1825} carries two linked variants (R53W missense and D125D
synonymous), suggesting a locally elevated substitution rate or a selective
sweep at this locus.

\begin{table}[h]
\centering
\caption{Twelve synapomorphies defining the clade~2\,+\,L6/L9/L10 node.
All are absent from clade~1, L5, L1--L4, L7, and \textit{M.~canettii}.}
\label{tab:synapo_c2_l6l9l10}
\begin{tabular}{@{}lllll@{}}
\toprule
SPDI & Gene & Locus & Effect & AA change \\
\midrule
NC\_000962.3:344956:A:G   & \textit{eccB3}  & Rv0283  & missense   & I312M \\
NC\_000962.3:445736:C:T   & Rv0368c         & Rv0368c & synonymous & P263P \\
NC\_000962.3:1699447:C:A  & Rv1508c         & Rv1508c & missense   & W149C \\
NC\_000962.3:2071191:C:T  & \textit{Rv1825} & Rv1825  & missense   & R53W \\
NC\_000962.3:2071409:C:T  & \textit{Rv1825} & Rv1825  & synonymous & D125D \\
NC\_000962.3:2150753:A:G  & Rv1903          & Rv1903  & missense   & T131A \\
NC\_000962.3:2453932:C:A  & Rv2191          & Rv2191  & missense   & R39S \\
NC\_000962.3:2501667:C:T  & Rv2228c         & Rv2228c & synonymous & L357L \\
NC\_000962.3:2789236:T:C  & \textit{plsB2}  & Rv2482c & missense   & D16G \\
NC\_000962.3:3393310:C:T  & ---             & ---     & intergenic & --- \\
NC\_000962.3:3981328:G:A  & Rv3542c         & Rv3542c & missense   & H218Y \\
NC\_000962.3:4225400:C:T  & Rv3779          & Rv3779  & synonymous & T139T \\
\bottomrule
\end{tabular}
\end{table}

The divergence between the two sister groups is strikingly asymmetric.
%
Clade~2 possesses 904~SPDIs absent from the L6/L9/L10 core, of which
23 are truly exclusive (Section~\ref{sec:results:stratification}).
%
Conversely, only 14~SPDIs are in the L6/L9/L10 core but absent from the
clade~2 core.
%
Closer examination reveals that 10 of these~14 are actually present in
three of the four clade~2 species and are excluded from the core only
because a single species has lost them: \Msur{} lacks three (including
\textit{Rv2818c}/Csm6~Q304*, a CRISPR-associated nonsense mutation),
\Mmun{} lacks three, and \Mchi{} lacks one (the \textit{stp} serine/threonine
phosphatase frameshift).
%
Only \textbf{4~SPDIs} are genuinely absent from all four clade~2 species:
one intergenic, \textit{lpqF}~F297L (a lipoprotein), and two synonymous
variants (\textit{glfT2}~V23V, \textit{eccCa1}~T121T).
%
This minimal L6/L9/L10-specific signature (4~SPDIs, only one missense)
suggests that secondary human re-adaptation required remarkably little
clade-specific genomic innovation beyond the shared ancestral repertoire.

\subsection{Phylogenetic clade signatures}
\label{sec:results:clade_sig}

To identify synapomorphic variants that define the major MTBC clades, we
computed four phylogenetic signatures matching the established topology
(Table~\ref{tab:clade_signatures}).
%
Signatures~A (L1--L7--L2--L3--L4 without L5) and~B (L1--L7--L2--L3--L4
+ L5) both yielded zero SPDIs: the strict intersection of L1 through L4
core genomes is empty because variant calling against the L4-derived H37Rv
reference means that L4-specific positions appear as wild-type in L4 strains
while L1/L2/L3-specific positions appear as wild-type in their respective
lineages, leaving no shared derived variants across all four.
%
Signature~C, the grand clade (L5~+ both animal clades~+ L6/L9/L10),
identifies 10~synapomorphic SPDIs present in all nine constituent lineages
(L5, L6, L9, L10, \textit{M.~bovis}, \textit{M.~suricattae},
\textit{M.~dassievlei}, \textit{M.~mungi}, \textit{M.~chimpanzae}) and
absent from L1--L4, L7 (exhaustively confirmed in all 81 L7 genomes),
and \textit{M.~canettii} (0/283 strains, except for two sporadic hits
at 1/283 each, consistent with sequencing noise).
%
Among the 10 grand-clade synapomorphies, three are frameshifts in
functionally important genes: \textit{ald}~(Rv2780; L-alanine
dehydrogenase, involved in nitrogen metabolism and dormancy),
\textit{cyp142}~(Rv3518c; cytochrome~P450 involved in host cholesterol
catabolism), and Rv1101c (hypothetical protein).
%
Four missense variants affect key processes: \textit{birA}~p.D211N
(Rv3279c; biotin ligase/transcriptional repressor), \textit{vapC48}~p.G34R
(Rv3697c; toxin--antitoxin persistence), Rv0966c~p.W186G (radical
aromatic$\to$aliphatic substitution), and Rv0986~p.V228L.
%
Signature~D identifies 192~SPDIs specific to L5 and absent from all other
lineages, confirming its long independent evolution since the basal node
of the grand clade.

\begin{table}[h]
\centering
\caption{Phylogenetic clade signatures.}
\label{tab:clade_signatures}
\begin{tabular}{@{}llr@{}}
\toprule
Signature & Definition & SPDIs \\
\midrule
A & L1--L7--L2--L3--L4 (excl.\ L5), absent from rest & 0 \\
B & L1--L7--L2--L3--L4 + L5, absent from rest & 0 \\
C & Grand clade (L5 + animals + L6/L9/L10), absent L1--L4/L7/\textit{canettii} & 10 \\
D & L5-specific, absent from all others & 192 \\
\bottomrule
\end{tabular}
\end{table}

\subsection{L5 as a genomic bridge}
\label{sec:results:l5}

L5 shares 216 SPDIs with the secondarily human-adapted clade (L6/L9/L10) that are absent
from strict human lineages, including notable LoF mutations in \textit{mce1R},
\textit{cyp142}, and \textit{tgs3}.
%
Conversely, L5 shares only two SPDIs with strict human lineages (L1--L4),
placing it genomically closer to the animal/basal grade than to the strict
human-adapted core.
%
The 10~grand-clade synapomorphies (Signature~C), shared by all lineages
on the L5/animal/L6--L9--L10 branch but absent from the strict human
clade (L1--L4--L7) and \textit{M.~canettii}, define the deepest genomic
split within MTBC after the root divergence.
%
The high proportion of functionally impactful variants (3~frameshifts and
4~non-conservative missense among 10~SPDIs) suggests that the ancestral
divergence event was accompanied by substantial metabolic rewiring,
particularly in nitrogen metabolism (\textit{ald}), cholesterol catabolism
(\textit{cyp142}), and biotin homeostasis (\textit{birA}).
%
The absence of any detectable synapomorphy for the strict human clade
(Signatures A and B both equal zero) reflects the inherent limitation of
SNP-based variant calling against the H37Rv reference, which belongs to L4.
This does not imply absence of derived characters in L1--L4--L7 but rather
that the SPDI framework is blind to positions where H37Rv itself carries
the derived allele.
%
The 192~L5-specific SPDIs (Signature~D) confirm the long independent
evolutionary trajectory of L5 within the grand clade, consistent with its
basal phylogenetic position and its restricted geographic distribution in
West Africa.

% ─────────────────────────────────────────────────────────────────────────────
\section{Discussion}
\label{sec:discussion}

\subsection{An evolutionary scenario for MTBC host-range specialization}
\label{sec:disc:evolution}

The genomic signatures identified here allow reconstruction of a comprehensive
evolutionary scenario for MTBC host-range specialization, spanning from a
generalist ancestor to modern specialist ecotypes.

The ancestral MTBC likely represented a broad-spectrum pathogen with moderate
genomic stability, as evidenced by the 388 shared core SPDIs present in both
animal and secondarily human-adapted lineages. This ancestral state maintained functional
redundancy in metabolic pathways and possessed the complete suite of virulence
systems (ESX-1 through ESX-5, intact PKS/PDIM machinery, and diverse PE/PPE
repertoires) without the host-specific modifications observed in modern lineages.

Ecological divergence initiated contrasting evolutionary trajectories.
Animal-adapted populations experienced massive activation of mobile genetic
elements, particularly IS6110, leading to a 25-fold increase in transposition
activity compared to human lineages. This ``genomic instability phase'' drove
convergent evolution across 333 genes in four independent animal lineages,
with IS6110 preferentially targeting PE/PPE loci (particularly \textit{PPE55})
to generate antigenic diversity. Metabolic streamlining occurred through
selective loss-of-function in peripheral pathways (vitamin B12 synthesis,
aldehyde metabolism) while preserving essential virulence functions.

The human trajectory involved a contrasting ``genomic stabilization'' event,
characterized by complete cessation of IS activity. The 10~grand-clade
synapomorphies, including frameshifts in \textit{ald} and \textit{cyp142},
mark the ancestral divergence of the L5/animal/L6--L9--L10 branch from
the strict human clade, while 14~L6/L9/L10-exclusive SPDIs
(including \textit{stp} frameshift and \textit{Csm6} Q304* nonsense)
represent the molecular signatures of secondary human re-adaptation.
L5 occupies a basal position within the grand clade, with 192~unique
SPDIs reflecting its long independent evolution, and likely represents
the transitional state between generalist and human-specialist forms.

The resulting dichotomy represents two fundamentally different evolutionary
strategies: animal lineages pursuing ``diversifying selection'' through ongoing
genomic instability, and human lineages achieving ``optimizing selection''
through stabilization and transmission enhancement. This model predicts that
zoonotic spillover risk remains low due to the extensive genomic specialization
in both ecotypes, and suggests that therapeutic targeting should account for
the distinct evolutionary pressures shaping each host-adapted clade.

This study reveals several findings that extend beyond previous MTBC evolutionary
studies (Table~\ref{tab:novel_findings}). While IS6110 transposition dynamics have been reported across MTBC
lineages \cite{Gonzalo2018}, the 25-fold differential activity between animal
and human ecotypes, and the unprecedented targeting of a single PE/PPE locus
(\textit{PPE55}) representing 68.4\% of total IS activity, have not been
previously described. Similarly, while convergent evolution has been documented
in bacterial pathogens (typically 50--200 genes in \textit{Pseudomonas aeruginosa}
\cite{Marvig2015}), the 333-gene convergent set identified here represents an
exceptional 8.3\% of the MTBC genome, suggesting unusually strong selective
pressures during animal host specialization. The 10~grand-clade synapomorphies define the deepest split within MTBC
after the root, and the 14~L6/L9/L10-exclusive SPDIs provide the first
SNP-level molecular markers of secondary human adaptation.

\subsection{Complex phylogenetic topology reveals multiple host-jumping events}
\label{sec:disc:complex_phylogeny}

Comprehensive phylogenetic analysis of thousands of MTBC genomes reveals a
complex evolutionary topology that contradicts the traditional model of simple
animal-human bifurcation. The data support a multi-host jumping scenario
where L6, L9, and L10 represent secondarily human-adapted lineages that evolved
from an animal-adapted intermediate rather than representing basal human lineages.
%
This interpretation is strongly supported by recent phylogenetic discoveries,
including the identification of L10 as a ``missing link'' in \textit{M.~africanum}
evolutionary history \cite{Guyeux2024L10} and the documentation of paraphyletic
relationships among animal-adapted MTBC members \cite{Brites2018}, where some
animal lineages are more closely related to certain human lineages than to
other animal lineages.

The proposed evolutionary scenario involves at least four independent
host-switching events: (1) initial divergence from a generalist African
ancestor, (2) specialization toward animal hosts in the primarily
animal-adapted clade, (3) secondary adaptation back to human transmission
giving rise to L6/L9/L10, and (4) independent human specialization in the
L1--L4 lineages. This model explains several previously puzzling observations:
the geographic co-localization of L5 and L6/L9 in West Africa despite their
phylogenetic separation from strict human lineages, the molecular signatures
of convergent adaptation toward human transmission observed in both L6/L9/L10
and L1--L4, and the apparent asymmetric burden of animal-differentiating variants
(904 SPDIs, of which 23 truly exclusive) versus secondarily human-adapted
exclusive variants (14 SPDIs), although rarefaction analysis shows this
asymmetry is largely driven by differential sampling depth
(Section~\ref{sec:disc:limits}).

The genomic evidence supporting this secondary adaptation hypothesis includes:
(1) the 10~grand-clade synapomorphies shared by L5, both animal clades, and
L6/L9/L10 but absent from L1--L4/L7 and \textit{M.~canettii}, confirming
that L6/L9/L10 belong to the same major branch as the animal lineages;
(2) the 14~L6/L9/L10-exclusive SPDIs marking the genomic modifications
acquired during re-adaptation to human transmission; and (3) the transitional
genomic composition of these lineages showing intermediate characteristics
between animal and strict human adaptations.
This complex phylogenetic framework fundamentally revises our understanding
of MTBC evolution and highlights the remarkable plasticity of this pathogen
in adapting to diverse mammalian hosts through recurrent host-jumping events.

\subsection{Epidemiological implications of secondary human adaptation}
\label{sec:disc:epidemiological_implications}

The recognition that L6/L9/L10 represent secondarily human-adapted lineages
has potential implications for MTBC epidemiology. The evolutionary legacy of
their animal-adapted intermediate phase could conceivably confer distinct
transmission dynamics, tissue tropism, or immune-evasion profiles compared to
strict human lineages (L1--L4), and the demonstrated capacity for
bidirectional host-jumping challenges traditional unidirectional spillover
models. However, the present study provides genomic evidence only; whether
these genomic differences translate into measurable epidemiological or
clinical differences remains an open question that will require dedicated
clinical and epidemiological studies in regions where L6/L9/L10 co-circulate
with animal-adapted MTBC members.

\subsection{Phylogeographic evidence for L5 as an evolutionary bridge}
\label{sec:disc:l5_bridge}

The intermediate phylogenetic position of L5 between animal-adapted and strictly
human-adapted lineages, combined with its West African geographic restriction
\cite{Gehre2013,YeboahManu2017}, suggests this lineage represents an evolutionary
bridge in MTBC host specialization. L5 shares 216 SPDIs with secondarily human-adapted lineages
(L6/L9/L10) that are absent from strict human lineages (L1--L4), and
all 10~grand-clade synapomorphies including the \textit{ald} and
\textit{cyp142} frameshifts.
Additionally, L5 harbours 192~unique SPDIs absent from every other lineage,
including loss-of-function mutations in \textit{mce1R} (mammalian cell entry
regulator) and lipid metabolism genes (\textit{fadD32}), that suggest
adaptations to intracellular survival distinct from both animal and strict
human lineages.

West Africa represents the only region where six of seven major MTBC lineages
are endemic \cite{YeboahManu2017}, with L5 showing maximum genetic diversity
in this region \cite{Gehre2013}. This phylogeographic pattern supports a
``West African Bridge'' hypothesis wherein L5 represents a relictual population
from an ancient spillover event that preceded the evolution and global expansion
of strict human-adapted lineages. The shared mutations between L5 and
secondarily human-adapted lineages (L6/L9/L10) likely reflect adaptations
acquired at the
ancestral node of African human-adapted MTBC, while the subsequent evolution
of L1--L4 involved additional specializations for enhanced human-to-human
transmission that are absent from both L5 and animal-adapted lineages.

\subsection{Temporal calibration of MTBC evolution}
\label{sec:disc:temporal}

Published molecular dating estimates place the MTBC root at
$\sim$70{,}000~years ago and the animal--human divergence at
$\sim$50{,}000~years ago \cite{Comas2013}, predating the major demographic
transitions associated with agriculture. Within this framework, the
10~grand-clade synapomorphies mark the deep split shared by L5, animal
lineages, and L6/L9/L10, while the 14~L6/L9/L10-exclusive SPDIs likely
arose during the subsequent secondary human adaptation. Precise dating of
intermediate nodes remains uncertain and will require dedicated Bayesian
phylogenomic analyses with broader taxon sampling.

\subsection{Speculative evolutionary mechanisms for secondary human adaptation}
\label{sec:disc:reemergence_mechanisms}

The secondary adaptation of L6/L9/L10 from animal-adapted intermediates back
to human hosts is one of the most intriguing features of MTBC evolution, but
the mechanisms remain speculative. One plausible hypothesis is \emph{latent
pre-adaptation}: animal-adapted lineages may have retained vestigial machinery
for human transmission that could be reactivated with relatively few genomic
changes, consistent with the small number of L6/L9/L10-exclusive SPDIs
identified here. Climate-driven host compression during the Last Glacial
Maximum could have provided the ecological opportunity for such a switch, as
documented for other bacterial host-jumps \cite{Spyrou2022}. However, these
scenarios require dedicated experimental validation---particularly functional
testing of the 4~truly L6/L9/L10-exclusive variants, the 12~node
synapomorphies, and the grand-clade synapomorphies---before they can be
elevated beyond informed speculation.

\begin{table}[h]
\centering
\caption{Novel findings and their relationship to existing knowledge}
\label{tab:novel_findings}
\begin{tabular}{@{}p{4cm}p{4cm}p{4cm}p{2cm}@{}}
\toprule
\textbf{Discovery} & \textbf{Previous knowledge} & \textbf{Our contribution} & \textbf{Originality} \\
\midrule
IS6110 activity 25× higher in animal lineages & IS6110 distribution varies between lineages \cite{Gonzalo2018} & Quantified differential activity; mechanistic explanation & High \\
PPE55 IS6110 hotspot (68.4\% activity) & No prior reports of IS6110-PE/PPE targeting & First description of IS insertion hotspot in PE/PPE genes & Complete \\
333 convergent genes (8.3\% genome) & Convergent evolution: ~50-200 genes typical \cite{Marvig2015} & Largest bacterial pathogen convergent gene set described & High \\
10 grand-clade synapomorphies & Clade-defining SNPs not systematically catalogued & First SNP-level synapomorphy of the L5/animal/L6--L9--L10 branch & High \\
23 truly animal-exclusive SPDIs & Animal core defined by RD deletions only \cite{Brosch2002} & First ancestry-stratified SNP synapomorphy of the animal clade & High \\
0 convergent genes between animal clades & Animal clades assumed to share common adaptations & Phenotypic convergence via non-overlapping genomic routes & Complete \\
12 synapomorphies of C2\,+\,L6/L9/L10 node & No SNP-level markers for this branching point & First definition including \textit{eccB3} (ESX-3) and \textit{plsB2} (lipids) & High \\
4 truly L6/L9/L10-exclusive SPDIs & Secondary human adaptation assumed to require many changes & Minimal genomic innovation for host re-switching & High \\
\bottomrule
\end{tabular}
\end{table}

\subsection{Characterising the animal core genome: ancestry stratification}
\label{sec:disc:identity}

The 904 animal-differentiating SPDIs reported here are fixed in all four animal
lineages and absent from the L6/L9/L10 core, providing a comprehensive
characterisation of the animal core genome relative to its closest
human-adapted sister clade.
%
Ancestry stratification reveals that 68\% of these variants are ancestral
MTBC polymorphisms (present in \textit{M.~canettii}), reflecting the
well-known H37Rv reference bias rather than animal-specific innovation.
%
After filtering against all available MTBC lineages ($>$50 L4 sub-lineages,
L1--L3, L5, L7, and \textit{M.~canettii}), only 23~SPDIs (2.5\%) are
truly exclusive to the animal clade --- the genuine synapomorphies of the
animal-adapted MTBC ecotype.
%
Their coherence across lineages as divergent as \Msur{} and \Mchi{} --- estimated
to have shared a common ancestor over a million years ago \cite{Boritsch2016} ---
implies that these 23 mutations were already fixed in the ancestral animal clade
population, prior to the ecological divergence into individual host species.

This SNP-resolution definition complements and extends the established
Region of Difference (RD) framework for MTBC classification \cite{Brosch2002}.
%
RD deletions (notably RD7, RD8, RD9, and RD10) have long served as the
primary markers distinguishing animal-adapted from human-adapted MTBC
\cite{Brosch2002, Mostowy2005}, but they capture only large chromosomal deletions
and provide no information on the fine-grained functional divergence that
accumulated alongside them.
%
The 904-SPDI animal core includes 439 missense (48.5\%) and 40
loss-of-function variants distributed across 611 loci, revealing a layer of
molecular variation invisible to deletion-based typing.
%
The elevated dN/dS proxy of the animal core (1.65) should be interpreted
with caution, as the majority of these variants are ancestral MTBC
polymorphisms rather than animal-specific innovations; the signal likely
reflects the overall composition of ancient MTBC variation at positions
where H37Rv carries derived alleles.

The functional composition of the core is instructive.
%
While ``conserved hypotheticals'' account for the largest absolute count
(265 SPDIs, 29.3\%), consistent with the baseline polymorphism rate of
poorly characterised genes, the next most represented categories are
intermediary metabolism \& respiration (146 SPDIs, 16.1\%), cell wall
\& cell processes (141 SPDIs, 15.6\%), PE/PPE (60 SPDIs, 6.6\%), and lipid
metabolism (53 SPDIs, 5.9\%).
%
This profile differs markedly from the secondarily human-adapted clade, which harbours
only 14 lineage-exclusive SPDIs at full sampling depth ($n = 1{,}132$) ---
nine of which are absent from any animal lineage --- although rarefaction
analysis demonstrates that this asymmetry is largely attributable to the
44-fold difference in sampling depth between the two clades
(Section~\ref{sec:disc:limits}).
%
Taken together with the 388 SPDIs shared by both animal and secondarily human-adapted clades
(the ancestral MTBC core), these data paint a picture in which the animal clade
underwent a burst of adaptive fixation relative to the ancestral state, while
the secondarily human-adapted lineages remained comparatively close to the ancestral genomic
configuration.
%
This is consistent with the hypothesis that the MTBC ancestor was a generalist
pathogen of African mammals, and that the animal-adapted lineages underwent
accelerated molecular evolution as they specialised to increasingly restricted
host niches \cite{Comas2013, Boritsch2016}.

\subsection{PKS/PDIM remodelling as a hallmark of animal host adaptation}
\label{sec:disc:pks}

The concentration of adaptive divergence in the PKS/PDIM pathway is consistent
with the established role of PDIM in modulating host immune recognition.
%
PDIM shields the mycobacterial surface from Toll-like receptor detection and
facilitates phagosomal escape \cite{Camacho2001, Rousseau2004}, and its
structural modification could alter the nature of the host--pathogen interaction
in a host-specific manner.
%
Recent work demonstrates that metabolic environment directly governs PDIM
maintenance: propionate availability prevents loss of PDIM production
\cite{Mulholland2024}, providing a mechanistic link between host-specific
nutrient environments and the selective pressures on PKS/PDIM genes observed
here. Furthermore, PDIM loss has been shown to differentially alter ESX-1
substrate secretion \cite{vanWijk2024}, suggesting that the coordinated
divergence of PKS/PDIM and ESX systems in animal lineages may be
functionally coupled rather than independent.
%
The absence of LoF in the central ppsA--E/mas/mmpL7 pathway suggests that PDIM
production itself remains essential across all MTBC ecotypes, while auxiliary
PKS enzymes modulating PDIM chain length or accessory lipid composition are
subject to relaxed or positive selection.

\subsection{PE/PPE surface variability and host-specific immune evasion}
\label{sec:disc:pppe}

The PE/PPE gene families constitute one of the most distinctive features of the
MTBC genome, accounting for approximately 10\% of its coding capacity.
%
These proteins, characterised by conserved N-terminal PE or PPE domains and
hypervariable C-terminal regions, are predominantly surface-exposed or secreted,
and have been implicated in antigenic variation, modulation of host immune
responses, and manipulation of phagosomal trafficking
\cite{Banu2002, Sampson2011, Peng2022}.
%
Our data reveal that this gene family is disproportionately targeted by
animal-clade mutations: 60 of the 904 animal-differentiating SPDIs (6.6\%)
fall within PE/PPE genes, spanning 30 distinct loci, a proportion significantly
enriched relative to the ancestral core ($p = 0.0002$, Fisher's exact test)
and to the secondarily human-adapted exclusive set, which contains no PE/PPE SPDIs whatsoever.

Among the 60 PE/PPE SPDIs, 28 are missense variants and six are loss-of-function
mutations (frameshifts in \textit{PE\_PGRS47}, \textit{PPE5}, \textit{PPE6},
\textit{PPE7}, \textit{PPE55}, and \textit{PPE58}), suggesting a dual dynamic:
ongoing diversification of surface-exposed residues on one hand, and outright
pseudogenisation of specific family members on the other.
%
This pattern is consistent with the well-described birth-and-death evolution
of PE/PPE genes, whereby individual members can be silenced while the overall
family retains functional diversity \cite{Veyrier2009}.

\textit{PPE55} (Rv3347c) emerges as the most intensely mutated locus in the
entire animal-differentiating core, with 13 SPDIs across all four lineages.
%
These include a frameshift shared by all four animal species --- fixed in the
ancestral animal clade --- alongside at least eight non-synonymous substitutions
distributed along the C-terminal hypervariable region, several of which are also
shared across all four lineages (L998F, M999I, N992T, L2259P, F1003D).
%
The co-occurrence of a frameshift with maintained missense variation in the same
gene is unusual and may reflect a functional shift: the frameshift could
truncate a surface-exposed domain while preserving a structurally distinct
isoform still subject to immune-driven diversification.

The signal of PE/PPE divergence is further amplified by the convergence analysis.
%
PPE54, PPE5, PPE6, and PPE8 each carry independent, non-identical mutations in all
four animal lineages --- the strongest possible evidence for parallel adaptive
pressure on these loci.
%
The convergent mutations in PPE54 and PPE5 include inframe indels in their
C-terminal domains, a type of variation associated with rapid antigenic
diversification under immune selection \cite{Fishbain2015}.

A mechanistically coherent picture emerges when this PE/PPE remodelling is
considered alongside the ESX-5 secretion system data.
%
ESX-5 is the dominant type-VII secretion system responsible for translocating
PE/PPE proteins to the mycobacterial surface and into the host cell
\cite{Houben2012}, and it is precisely ESX-5 that carries the highest number
of animal-differentiating SPDIs among the five ESX systems (nine SPDIs, all missense
or synonymous; Table~\ref{tab:esx}).
%
We propose that the coordinated divergence of ESX-5 and its PE/PPE substrate
repertoire reflects co-evolution of a secretion apparatus with its cargo,
optimising the presentation of host-specific surface antigens without
compromising the integrity of the secretion pathway itself --- an interpretation
supported by the complete absence of loss-of-function mutations in any ESX system
across all four animal lineages.

\subsection{ESX secretion systems: conserved function, altered substrate specificity}
\label{sec:disc:esx}

Perhaps the most striking finding of the ESX analysis is its negative result: across
all 24 animal-differentiating SPDIs distributed over the five type-VII secretion systems,
not a single loss-of-function mutation is observed.
%
This conservation is biologically expected --- ESX-1 is required for phagosomal
disruption and full virulence \cite{Hsu2003, Derbyshire2012}, ESX-3 is essential
for iron and zinc acquisition and cannot be deleted without loss of viability
\cite{Serafini2009, Siegrist2009}, and ESX-5 is the primary translocase for the
PE/PPE surface protein repertoire \cite{Houben2012} --- yet it is not trivially
guaranteed, given the extensive pseudogenisation observed in other functional
systems of the same genomes.
%
Its confirmation therefore constitutes meaningful evidence that type-VII secretion
represents a conserved architectural constraint across all MTBC ecotypes,
regardless of host range.

The pattern of missense divergence, however, is system-specific and biologically
interpretable.
%
ESX-5 accumulates the most animal-differentiating variation (nine SPDIs), consistent
with its role as the secretion conduit for the PE/PPE proteins that are themselves
the most diverged gene family in this dataset (Section~\ref{sec:disc:pppe}).
%
The affected components include the channel subunit EccB5 (p.R41H), the membrane
component EccD5 (p.G261D), and the secreted substrates EsxL (p.A12D) and EsxO
(p.G22A), suggesting that both the translocon and its cargo have co-evolved.
%
Notably, the EsxL A12D substitution falls within the conserved WXG motif of
ESAT-6 family proteins, a region critical for dimerisation and membrane
interaction \cite{Poulsen2014}, and its fixation in the animal clade core implies
functional significance.

ESX-1, the primary virulence secretion system, carries five animal-differentiating SPDIs,
but these are restricted to accessory Esp proteins (EspA p.T192I; EspK p.C729S and
p.S194F) rather than the core ATPases or the primary secreted effectors EsxA and
EsxB.
%
EspA and EspK are regulatory and stabilising components of the ESX-1 secretion
complex \cite{Chen2013}, and their divergence in the animal clade may reflect
fine-tuning of effector secretion kinetics or specificity in animal host
phagosomes, rather than a qualitative change in secretion competence.

ESX-3 presents a subtler but biologically consequential pattern: the missense
variant EccC3 D1146N (shared by all four animal lineages) affects the Walker~B
motif of the core ATPase that drives substrate translocation \cite{Poulsen2014}.
%
ESX-3 couples iron and zinc uptake to the recycling of mycobactin siderophores
\cite{Serafini2009}, and its modification in the animal clade may reflect
adaptation to the distinct metal availability profiles of the tissues colonised
in their respective animal hosts.
%
This hypothesis is further supported by the human lineage trajectory: EccE3
N217D, a missense in a different ESX-3 subunit, is found across human lineages
(L1--L4, L5, L6/L9/L10) and absent from the animal clade, although it also
represents an ancestral MTBC variant relative to H37Rv rather than a
human-specific innovation.
%
Animal and human ecotypes therefore appear to have independently remodelled
ESX-3 at different structural positions, converging on a shared outcome ---
altered iron/zinc acquisition --- through divergent molecular means.

Taken together, these data support a model in which the five ESX systems are
subject to strong purifying selection at the pathway level but to ongoing
adaptive evolution at the level of individual substrate-interacting and
regulatory components.
%
The animal and human ecotypes have each accumulated system-specific changes
--- predominantly in ESX-5 and its PE/PPE cargo for the animal clade, and in
ESX-1 and ESX-3 scaffold components for human lineages --- reflecting the
distinct immunological environments imposed by their respective hosts.
%
This model of ``tuning without destruction'' is further supported by the
recent discovery that homopolymeric tracts near ESX-1 regulators generate
phase variants that modulate secretion levels and host fitness
\cite{Hartman2025}, providing an additional, non-mutational mechanism for
ESX functional adjustment.

\subsection{Convergent evolution implicates lipid metabolism as an adaptive axis}
\label{sec:disc:convergence}

After subtracting the 904 mutations fixed in the animal core, 333~genes
carry independent, non-identical mutations in all four animal lineages.
%
The dominant functional signal among these convergently evolving loci is
lipid metabolism: at least 18 polyketide synthase and fatty-acid metabolism genes
--- including \textit{pks2}, \textit{pks3}, \textit{pks4}, \textit{pks6},
\textit{pks8}, \textit{pks9}, \textit{pks12}, \textit{pks13}, \textit{ppsA--C},
\textit{ppsE}, \textit{fadH}, \textit{fadD5}, \textit{fadD15}, and \textit{plsB2}
--- are convergently mutated in all four lineages.
%
This result survives the conservative robustness test restricted to the two
multi-genome lineages (\Msur{} and \Mdas{}), confirming that the signal is not
an artefact of single-genome sampling.

Two non-exclusive mechanisms can account for convergent mutation accumulation
in the same genes across independently evolving lineages: parallel positive
selection driven by a common adaptive pressure, or independent relaxation of
purifying selection in genes whose function has become partially dispensable
in a new ecological context.
%
Our data support both, operating on different gene subsets.
%
The elevated missense/synonymous ratios observed for \textit{ppsA} (2.0),
\textit{pks8} (2.0), and \textit{pks13} (2.0) are suggestive of positive or
relaxed selection, although these proxy values are derived from small absolute
counts and formal $\omega$ estimation (e.g.\ PAML site models) would be
required to confirm this interpretation.
%
In contrast, the convergent frameshift in \textit{pks6} (independently acquired in
\Msur{}) and the frameshift in \textit{fadD15} (present in all four lineages)
point to relaxed purifying selection: when a gene becomes dispensable, independent
loss-of-function events become the expected outcome under drift.

The biological coherence of this dichotomy is informative.
%
The PDIM central pathway (ppsA--E, mas, mmpL7) retains zero loss-of-function
mutations across all lineages --- indicating that PDIM production itself is
non-negotiable --- while the auxiliary PKS enzymes that modulate the accessory
lipid repertoire are progressively shed or diversified in a lineage-specific
manner.
%
All four animal hosts share one biological constraint: they impose a physically
distinct tissue niche (meerkats: pulmonary/lymphatic; mongooses: nasal passages;
hyraxes: lymph nodes; chimpanzees: respiratory tract) with different macrophage
lipid environments \cite{Boritsch2016}, potentially selecting for parallel
remodelling of the mycobacterial surface lipid profile.

Beyond the PKS/PDIM axis, the convergence analysis identifies two regulatory
genes mutated independently in all four animal lineages: \textit{pstP}, encoding
a phosphoserine/threonine phosphatase involved in cell-wall remodelling
\cite{Chopra2003}, and \textit{tcrY}, a two-component response regulator.
%
Their convergent missense divergence suggests that signal transduction pathways
governing cell-wall dynamics have also been independently retuned across the
animal clade, possibly in response to differences in the host phagosomal
environment.
%
These convergent patterns at the SNP level complement recent pan-genome analyses
showing that gene content diversification in MTBC is shaped by both phylogenetic
and ecological signatures, with each ecotype carrying its own accessory genome
\cite{Zwyer2024}, and that convergent deletions in PE/PPE and other regions
contribute to the small, closed MTBC pangenome \cite{Orgeur2025}.

\subsection{Phenotypic convergence despite genomic divergence between animal clades}
\label{sec:disc:inter_clade}

The comparison between the two animal clades provides a striking illustration
of how convergent ecological outcomes can arise through entirely different
molecular routes.
%
Clade~1 (Bovis/Caprae/Orygis/La4/Microti/Pinnipedii) and clade~2
(Suricattae/Dassie/Mungi/Chimpanzé) share 413~core SPDIs, but this shared
repertoire consists entirely of grand-clade ancestral variants (all also present
in at least one L6/L9/L10 strain): not a single SPDI is shared by both animal
clades while being absent from all non-animal lineages.
%
In other words, the two animal clades have no common genomic ``animal identity''
beyond their shared phylogenetic background.

Their respective synapomorphies (2 for clade~1, 23 for clade~2) target
entirely non-overlapping gene sets.
%
Clade~1 retains only two truly exclusive SPDIs after filtering: one intergenic
and a missense variant in the antitoxin \textit{vapB15}~(N63D), suggesting a
minimal clade-specific genomic signature despite encompassing six species.
%
Clade~2, in contrast, shows 23~innovations including type-VII secretion
(\textit{eccC3}~D1146N in the ESX-3 ATPase), core transcriptional machinery
(\textit{rpoA}~V325F), and host cholesterol metabolism (\textit{hsaB}~A168G).
%
The absence of any convergently targeted gene between the two exclusive sets
is notable given that both clades face analogous selective pressures as
animal-restricted pathogens, and contrasts sharply with the 333-gene
convergent signal observed \emph{within} clade~2 across its four species.

This pattern has two important implications.
%
First, it suggests that the selective pressures driving animal host adaptation
operate at the pathway level (lipid remodelling, immune evasion, intracellular
survival) rather than at the level of individual genes, allowing multiple
molecular solutions to the same ecological challenge.
%
Second, it reinforces the view that ``animal-adapted MTBC'' is not a monolithic
ecotype but rather a paraphyletic assemblage of independently evolving lineages
whose convergence is phenotypic --- restricted host range, reduced human
transmissibility --- rather than genomic.
%
The striking asymmetry in clade-specific synapomorphies (2 for clade~1 vs~23
for clade~2), despite substantially deeper sampling of clade~1
(529 vs~26~genomes), suggests that the six-species clade~1 core is heavily
constrained: with more species in the intersection, most derived variants
are diluted away, leaving only the deepest-rooting mutations.
%
Clade~2, with fewer but more divergent species, retains more of its
branch-defining variants in the core.

\subsection{The clade~2\,/\,L6--L9--L10 sister relationship: divergence from a shared node}
\label{sec:disc:c2_l6l9l10}

The 12~synapomorphies defining the clade~2\,+\,L6/L9/L10 node
(Table~\ref{tab:synapo_c2_l6l9l10}) provide the first SNP-level definition
of this phylogenetic branching point.
%
The presence of \textit{eccB3}~I312M (ESX-3 secretion) and
\textit{plsB2}~D16G (membrane lipid biosynthesis) among these 12 variants
echoes the recurrent targeting of secretion and lipid pathways across MTBC
diversification, and suggests that modifications in these systems were already
underway before the split between animal-adapted and secondarily human-adapted
lineages.

The most striking result, however, is the near-total absence of a
L6/L9/L10-specific genomic signature.
%
Only 4~SPDIs are genuinely exclusive to L6/L9/L10 (absent from all four
clade~2 species), and only one encodes a missense change
(\textit{lpqF}~F297L).
%
The remaining 10~SPDIs that appear ``L6/L9/L10-exclusive'' when compared
against the clade~2 core are in fact present in three of the four clade~2
species: they are excluded from the core because a single species has
independently lost each variant --- \Msur{} (3~losses, including
\textit{Rv2818c}/Csm6~Q304*), \Mmun{} (3~losses), and \Mchi{}
(1~loss: the \textit{stp} phosphatase frameshift).
%
This observation underscores a methodological caveat: core-genome
intersections are sensitive to single outlier species, and
``group-exclusive'' sets should always be validated at the per-species level.

The asymmetry in divergence is remarkable: 904~SPDIs separate clade~2 from
the L6/L9/L10 core, but only 14 (of which 4~genuine) separate L6/L9/L10
from clade~2.
%
This implies that the secondary re-adaptation of L6/L9/L10 to human hosts
involved minimal clade-specific genomic innovation.
%
Rather than acquiring a new ``human-adapted'' gene set, L6/L9/L10 lineages
may have simply retained the ancestral repertoire inherited from the
clade~2\,+\,L6/L9/L10 common ancestor while clade~2 diverged extensively
during animal adaptation.
%
In this view, the ``return to human'' of L6/L9/L10 is less a re-adaptation
event and more a case of phylogenetic inertia: these lineages never fully
lost the capacity for human transmission, while clade~2 accumulated the
mutations that restricted it to animal hosts.

\subsection{Grand-clade synapomorphies and the reference-bias limitation}
\label{sec:disc:clade_sig}

The phylogenetic clade signature analysis reveals a striking asymmetry: the
grand clade (L5 + animal lineages + L6/L9/L10) is defined by 10~robust
synapomorphies, while the strict human clade (L1--L4--L7) yields none.
%
This asymmetry is not biological but methodological: SNP calling against
the L4-derived H37Rv reference genome means that positions where L4 carries
a derived allele appear as wild-type and are invisible to variant detection.
%
A future analysis using an outgroup reference (e.g., \textit{M.~canettii})
or an imputed ancestral reference \cite{Malaga2024} would likely reveal a
symmetric set of derived characters on the strict-human branch.

The 10~grand-clade synapomorphies are functionally striking, with 70\%
(3~frameshifts + 4~non-conservative missense) predicted to affect protein
function.
%
The \textit{ald} (Rv2780) frameshift disrupts L-alanine dehydrogenase, a
key enzyme in nitrogen metabolism whose activity is critical during
dormancy and the transition between active replication and non-replicating
persistence states \cite{Hutter2004}.
%
The \textit{cyp142} (Rv3518c) frameshift affects a cytochrome~P450 involved
in host cholesterol catabolism, a pathway essential for intracellular survival
\cite{Johnston2010}; its disruption at the base of the grand clade
suggests an early metabolic shift away from cholesterol-dependent energy
production.
%
\textit{birA} p.D211N (Rv3279c) modifies the biotin protein
ligase/transcriptional repressor, potentially altering biotin homeostasis
and its downstream regulation of lipid biosynthesis \cite{Dey2010}.
%
\textit{vapC48} p.G34R (Rv3697c) affects a toxin component of the
toxin--antitoxin system implicated in persistence and dormancy regulation.

The 14~L6/L9/L10-exclusive SPDIs represent the molecular signatures
acquired during secondary human re-adaptation.
%
Among these, the \textit{stp} (Rv2333c) frameshift disrupts a
serine/threonine phosphatase that regulates cell-division and virulence
signalling \cite{Bach2008}, and \textit{Csm6} (Rv2818c) p.Q304* truncates
a CRISPR-associated RNase, potentially altering the CRISPR defence system
of these lineages.
%
The small number of L6/L9/L10-exclusive variants (14 SPDIs) compared to
the animal-differentiating core (904 SPDIs) reflects the relatively recent
secondary adaptation event and the minimal genomic modifications required
for the animal-to-human host transition.

The 192~L5-specific SPDIs confirm the long independent evolution of L5
since its divergence from the rest of the grand clade.
%
This massive private variant repertoire, combined with its basal
phylogenetic position, its restricted geographic distribution in West
Africa \cite{Gehre2013}, and its sharing of all 10~grand-clade
synapomorphies, firmly anchors L5 as a relictual lineage that preserves
the ancestral state of the grand-clade branch while having accumulated
extensive lineage-specific modifications during its independent evolution.

\subsection{Limitations}
\label{sec:disc:limits}

A key limitation of the per-species analyses concerns the unequal sampling depth
across animal lineages.
%
While \Msur{} is represented by 17 genomes and \Mdas{} by 7, both \Mmun{} and
\Mchi{} are each represented by a single sequenced genome, yielding 2{,}190 and
2{,}486 SPDIs respectively.
%
For single-genome lineages, all variants of the sequenced strain are treated as
fixed, without the ability to distinguish truly fixed mutations from within-lineage
polymorphisms.
%
This inflates SPDI counts by an estimated factor of $\times1.11$ (Mungi) and
$\times1.26$ (Chimpanzé) relative to the multi-genome Suricattae core.
%
Crucially, 61.9\% of Mungi SPDIs and 53.4\% of Chimpanzé SPDIs are independently
confirmed in the Suricattae genome set, and a sensitivity analysis restricting the
animal group to \Msur{} and \Mdas{} retains 100\% (904/904) of the animal-differentiating
core SPDIs.
%
The major biological signals reported here --- PE/PPE enrichment, PKS/PDIM pathway
divergence, and absence of ESX LoF --- are therefore robust to the inclusion or
exclusion of single-genome lineages.
%
Species-level conclusions for \Mmun{} and \Mchi{}, including dN/dS proxy values and
lineage-specific LoF mutations (e.g.\ \textit{pks7}~p.Q1265* in \Mchi{}),
should be interpreted with caution pending additional genomes from these lineages.

A second, more fundamental limitation concerns the asymmetric sampling depth
between the secondarily human-adapted clade ($n = 1{,}132$) and the animal
clade ($n = 26$).
%
The core genome (intersection) shrinks monotonically with sample size: a
rarefaction analysis subsampling L6/L9/L10 strains to $n = 26$ (matching the
animal sampling depth) yields a core of $\sim$1{,}400~SPDIs (vs 402 at full
depth), with $\sim$279~L6/L9/L10-exclusive SPDIs (vs 14) and only
$\sim$175~animal-differentiating SPDIs (vs 910).
%
The asymmetry between animal-differentiating (904) and secondarily human-adapted
exclusive (14) SPDIs reported in the main analysis therefore reflects
\emph{sampling depth} rather than true biological asymmetry in evolutionary
tempo.
%
At equal sampling depth ($n = 26$ each), the ratio would approach
$\sim$175:279 in favour of the secondarily human-adapted clade, suggesting
that --- if anything --- L6/L9/L10 lineages harbour more private variation
than the animal clade when measured fairly.
%
The functional conclusions of this study (PE/PPE enrichment, PKS/PDIM
remodelling, ESX conservation, convergent evolution, clade synapomorphies) are
robust to this bias, as they rely on the \emph{identity} and functional
annotation of variants rather than on raw counts. Nevertheless, all
quantitative comparisons of exclusive SPDI counts between clades with
unequal sampling depths should be interpreted with this caveat in mind.

Despite the extensive genomic divergence documented here, cross-referencing
the animal-differentiating core against the WHO 2022 drug resistance catalogue
revealed no variant classified as resistance-associated (Grades~1--2).
All 29~SPDIs falling within canonical DR loci are synonymous, lineage
markers, or Grade~4 (not associated). This negative result is
practically relevant: it confirms that standard molecular drug
susceptibility testing targets remain applicable to animal-adapted
MTBC lineages, and that animal-clade polymorphisms such as
\textit{katG}~R463L or \textit{gyrA}~S95T should not be misinterpreted
as resistance markers in veterinary or One Health surveillance contexts.

% ─────────────────────────────────────────────────────────────────────────────
\section{Conclusion}
\label{sec:conclusion}

This study fundamentally revises our understanding of MTBC evolution through
the discovery of complex multi-host evolution involving secondary human
adaptation. Based on comprehensive phylogenetic analysis supported by
\num{274800} whole-genome sequences, we demonstrate that the traditional
model of simple animal-human bifurcation is incorrect. Instead, MTBC evolution
involves multiple host-jumping events, with lineages L6, L9, and L10 representing
secondarily human-adapted lineages that evolved from animal-adapted intermediates
rather than constituting basal human lineages as previously assumed.
%
This phylogenetic breakthrough, supported by convergent evidence from molecular
signatures, temporal calibration, and geographic distribution patterns,
establishes a new evolutionary framework where the primarily animal-adapted
clade (characterised by 904 differentiating variants, 23 truly exclusive) and the
secondarily human-adapted lineages (14 exclusive variants at full depth,
noting sampling-depth bias; Section~\ref{sec:disc:limits})
represent distinct evolutionary solutions to host specialization.
%
Among the functional systems most diverged from human MTBC, four emerge as
primary axes of host-range differentiation: the PE/PPE gene family, which is
significantly enriched for animal-differentiating mutations and convergently diversified
across all four independent lineages; the PKS/PDIM lipid biosynthesis pathway,
whose central phthiocerol-synthesis module is remodelled without loss of function
while auxiliary PKS enzymes are progressively shed; the type-VII ESX
secretion systems, which maintain strict functional integrity in both ecotypes
but undergo non-overlapping missense divergence, with ESX-5 and its PE/PPE cargo
preferentially targeted in the animal clade and ESX-1 and ESX-3 scaffold
components diverging in human lineages; and insertion sequence (IS) element
activity, which shows a 25-fold enrichment in animal lineages (2.1\% vs 0.0\%),
with IS6110 preferentially targeting the \textit{PPE55} locus and contributing
to ongoing genomic instability in animal-adapted ecotypes.

The convergent evolution analysis provides particularly strong evidence for the
role of lipid metabolism as an adaptive axis: 333 genes (8.3\% of the genome),
dominated by polyketide synthase family members, carry independent mutations
in all four animal lineages---representing the largest documented convergent
gene set in bacterial pathogen evolution. This exceptional scale of convergence,
with dN/dS proxies consistent with positive selection in
the PDIM core synthases alongside relaxed purifying selection in auxiliary PKS
enzymes.
%
The 10~grand-clade synapomorphies --- including frameshifts in \textit{ald}
and \textit{cyp142} --- provide robust phylogenetic markers defining the
deepest intra-MTBC split after the root divergence. The 14~L6/L9/L10-exclusive
SPDIs, including the \textit{stp} frameshift and \textit{Csm6} Q304*
nonsense, mark the molecular events of secondary human re-adaptation.
The absence of detectable synapomorphies for the strict human clade reflects
the inherent reference bias of H37Rv-based variant calling, a methodological
limitation that should be addressed in future studies using outgroup-rooted
ancestral-state reconstruction.
%
The inter-clade comparison reveals that the two animal clades share no
exclusive genomic identity beyond the grand-clade ancestral background:
their 2 and 23 truly exclusive synapomorphies target entirely
non-overlapping gene sets, demonstrating that animal host adaptation is
phenotypically convergent but genomically divergent.
%
Furthermore, the clade~2\,+\,L6/L9/L10 sister relationship is defined by
12~node synapomorphies (including \textit{eccB3} and \textit{plsB2}), while
L6/L9/L10 itself possesses only 4~truly exclusive SPDIs, suggesting that
secondary human re-adaptation required remarkably little genomic innovation.
%
The discovery of secondary human adaptation has profound implications for
tuberculosis epidemiology, surveillance, and control strategies. Secondarily
adapted lineages may retain hybrid transmission characteristics, exhibit
altered zoonotic potential, and require specialized therapeutic approaches
that account for their unique evolutionary history. The demonstrated capacity
for bidirectional host-jumping fundamentally challenges traditional
unidirectional spillover models and necessitates revised risk assessment
frameworks that consider the potential for complex multi-host transmission
cycles in contemporary epidemiological contexts.

Beyond its immediate implications for MTBC research, this study establishes
a new paradigm for understanding bacterial pathogen evolution that extends
far beyond tuberculosis. The demonstration that animal-adapted lineages can
successfully re-adapt to human transmission through convergent molecular
mechanisms provides a crucial framework for predicting and monitoring emerging
infectious diseases. The evolutionary plasticity revealed here, operating
across tens of thousands of years and involving multiple independent host-switching
events, highlights the remarkable capacity of bacterial pathogens to respond
to changing ecological pressures through recurrent adaptation and counter-adaptation.

Several avenues remain open.
%
The single-genome representation of \Mmun{} and \Mchi{} is the principal
limitation of the present study; while our robustness analysis confirms that
the core findings are not driven by these singleton sequences, the
lineage-specific signals for these two species --- including the \textit{pks7}
p.Q1265* nonsense in \Mchi{} and the highest dN/dS proxy values in \Mmun{}
--- require validation against additional genomes.
%
Functional characterisation of the key variants identified here, in particular
the EccC3 D1146N substitution in the ESX-3 ATPase, the multiple PPE55 missense
cluster, and the \textit{ald} and \textit{cyp142} frameshifts in the grand
clade, would substantially
advance our mechanistic understanding of how secretion systems and
surface-lipid remodelling underlie the host-range boundaries within the MTBC.
Recent mechanistic studies demonstrating that PDIM lipids inhibit autophagy
\cite{Mittal2024} and are required for starvation-induced antibiotic tolerance
\cite{Block2023} provide functional context for the PKS pathway modifications
observed here.
%
The novel discoveries reported here have immediate translational applications.
The 10~grand-clade synapomorphies, 12~clade~2\,+\,L6/L9/L10 node markers,
and 4~truly L6/L9/L10-exclusive SPDIs provide molecular markers for
phylogenetic clade assignment and ecotype differentiation in diagnostic
settings, while the \textit{PPE55} IS6110
hotspot offers insights into ongoing genomic instability patterns that could
inform epidemiological predictions. The exceptional scale of
convergent evolution (333 genes) suggests that animal-adapted MTBC lineages
experience unusually strong and consistent selective pressures, potentially
enabling prediction of adaptive trajectories in emerging animal strains.
%
More broadly, the analytical framework developed here --- combining core SPDI
intersections, pathway-level dN/dS proxies, IS element activity quantification,
and convergent-evolution decomposition at scale --- provides a reusable template for comparative
genomic studies of host adaptation in other obligate bacterial pathogens.

% ─────────────────────────────────────────────────────────────────────────────
\bibliographystyle{unsrtnat}
\bibliography{references}

% ─────────────────────────────────────────────────────────────────────────────
\appendix

\section{Supplementary Tables}
\label{sec:supp_tables}

\begin{table}[h!]
\centering
\caption{Group composition and genome counts.}
\label{tab:groups}
\begin{tabular}{llrl}
\toprule
Group & Lineage(s) & $n$ genomes & Role \\
\midrule
Animal clade     & \textit{M.~suricattae}    & 17   & Main animal group \\
                 & \textit{M.~dassievlei}    & 7    & \\
                 & \textit{M.~mungi}         & 1    & \\
                 & \textit{M.~chimpanzae}    & 1    & \\
Second.\ human-adapted      & L6                        & 1089 & Comparison \\
                 & L9                        & 40   & \\
                 & L10                       & 3    & \\
Strict human     & L1--L4, L7                & $5\times$sub-lin + 81 & Comparison \\
L5               & L5                        & 443  & Outgroup \\
\bottomrule
\end{tabular}
\end{table}

\begin{table}[h!]
\centering
\caption{Core SPDI summary by group.}
\label{tab:core_summary}
\begin{tabular}{lrr}
\toprule
Category & SPDIs & Genes \\
\midrule
Animal clade core               & 1292 & 565 \\
Animal-differentiating (vs second.\ human) & 904  & 486 \\
\quad of which truly animal-exclusive & 23 & 20 \\
\quad of which ancestral (in \textit{canettii}) & 616 & --- \\
Ancestral (animal $\cap$ second.\ human) & 388 & 293 \\
Clade~2\,+\,L6/L9/L10 node synapomorphies & 12 & 10 \\
Second.\ human-adapted exclusive & 14   & 14 \\
\quad of which truly L6/L9/L10-only & 4 & 4 \\
Second.\ human absent all animal & 9    & 9 \\
Grand-clade synapomorphies (sig.\ C) & 10 & 8 \\
L5-specific (sig.\ D)          & 192  & --- \\
\bottomrule
\end{tabular}
\end{table}

\begin{table}[h!]
\centering
\caption{Animal-differentiating SPDIs in ESX type-VII secretion systems.}
\label{tab:esx}
\begin{tabular}{llllll}
\toprule
System & Gene & Effect & AA change & Predicted role & LoF \\
\midrule
ESX-1 & espA  & missense  & T192I & Secreted effector & No \\
ESX-1 & espD  & synonymous & G20G  & Effector secretion & No \\
ESX-1 & espJ  & synonymous & G243G & Effector modulation & No \\
ESX-1 & espK  & missense  & C729S & Effector modulation & No \\
ESX-1 & espK  & missense  & S194F & Effector modulation & No \\
\midrule
ESX-2 & eccA2 & synonymous & A189A & ATPase component & No \\
ESX-2 & eccE2 & synonymous & Y291Y & Membrane component & No \\
ESX-2 & esxD  & missense  & T49A  & Secreted substrate  & No \\
\midrule
ESX-3 & eccA3 & missense  & E6A   & ATPase component & No \\
ESX-3 & eccA3 & synonymous & A384A & ATPase component & No \\
ESX-3 & eccC3 & missense  & D1146N & Core ATPase & No \\
ESX-3 & PPE4  & synonymous & A185A & PE/PPE substrate & No \\
\midrule
ESX-4 & eccC4 & missense  & S1082G & Core ATPase & No \\
ESX-4 & mycP4 & synonymous & I158I & Mycosin protease & No \\
ESX-4 & eccB4 & synonymous & G417G & Channel component & No \\
\midrule
ESX-5 & eccB5 & missense  & R41H  & Channel component & No \\
ESX-5 & eccB5 & synonymous & P499P & Channel component & No \\
ESX-5 & eccC5 & synonymous & S1204S & Core ATPase & No \\
ESX-5 & eccD5 & missense  & G261D & Membrane component & No \\
ESX-5 & esxK  & synonymous & N39N  & PE/PPE secretion & No \\
ESX-5 & esxL  & missense  & A12D  & PE/PPE secretion & No \\
ESX-5 & esxL  & synonymous & A68A  & PE/PPE secretion & No \\
ESX-5 & esxO  & missense  & G22A  & PE/PPE secretion & No \\
ESX-5 & esxO  & synonymous & A83A  & PE/PPE secretion & No \\
\bottomrule
\end{tabular}
\end{table}

\section{Main Figures}
\label{sec:main_figs}

\begin{figure}[h]
\centering
\includegraphics[width=0.9\textwidth]{../figures/phylo_tree_english.png}
\caption{Phylogenetic relationships of MTBC animal and human lineages. Simplified tree showing the major evolutionary split between animal-adapted lineages (red), secondarily human-adapted lineages (teal), and modern human lineages (blue). The tree illustrates the fundamental animal-human divergence that underlies differential genomic adaptation patterns.}
\label{fig:phylo}
\end{figure}

\begin{figure}[h]
\centering
\includegraphics[width=0.9\textwidth]{../figures/functional_enrichment_english.png}
\caption{Functional enrichment comparison across MTBC lineage groups. Bar plots showing gene counts by functional category for animal-differentiating (red), secondarily human-adapted exclusive (teal), and ancestral core (blue) variants. Animal lineages show significant enrichment for PE/PPE proteins and lipid metabolism (***p < 0.001).}
\label{fig:enrichment}
\end{figure}

\begin{figure}[h]
\centering
\includegraphics[width=0.9\textwidth]{../figures/lof_comparison.png}
\caption{Loss-of-function mutation comparison across lineage groups. Bar chart showing counts of nonsense, frameshift, and start-loss mutations. Animal lineages exhibit higher LoF burden, particularly in PE/PPE and metabolic genes.}
\label{fig:lof}
\end{figure}

\begin{figure}[h]
\centering
\includegraphics[width=0.9\textwidth]{../figures/pks_pdim_overview.png}
\caption{PKS/PDIM pathway overview showing differential remodeling in animal vs human lineages. (A) Pathway schematic, (B) Mutation distribution across pathway components, (C) dN/dS proxy values indicating selection pressure, (D) Species-specific LoF patterns.}
\label{fig:pks_overview}
\end{figure}

\begin{figure}[h]
\centering
\includegraphics[width=0.9\textwidth]{../figures/pks_pdim_pathway_map.png}
\caption{Detailed PKS/PDIM biosynthetic pathway map. Enzymatic steps from fatty acid precursors to final PDIM products, with mutation counts and effect types overlaid on each gene. Red indicates high mutation burden.}
\label{fig:pks_pathway}
\end{figure}

\begin{figure}[h]
\centering
\includegraphics[width=0.9\textwidth]{../figures/convergence_english.png}
\caption{Convergent evolution scale and functional patterns in animal MTBC lineages. (A) Number of genes showing convergent mutations across 4, 3, or 2 animal species, representing 8.3\% of the MTBC genome. (B) Top functional categories of convergent genes, dominated by PKS lipid synthesis and PE/PPE proteins.}
\label{fig:convergence}
\end{figure}

\begin{figure}[h]
\centering
\includegraphics[width=0.9\textwidth]{../figures/is_activity_english.png}
\caption{Differential IS element activity between animal and human MTBC lineages. (A) IS activity comparison showing 25-fold higher insertion rate in animal lineages (2.1\% vs 0.0\%). (B) IS element type breakdown in animal lineages, with IS6110 dominance and PPE55 hotspot representing 68.4\% of total activity.}
\label{fig:is_activity}
\end{figure}

\begin{figure}[h]
\centering
\includegraphics[width=0.9\textwidth]{../figures/functional_categories_comparison.png}
\caption{Mycobrowser functional category distribution across SPDI groups. Stacked bar chart showing the proportion of each Mycobrowser functional category (intermediary metabolism \& respiration, cell wall \& cell processes, conserved hypotheticals, lipid metabolism, PE/PPE, etc.) for each comparative group: animal-differentiating (904), L6/L9/L10-exclusive (14), core ancestral (388), and other intersection sets. Note the elevated PE/PPE and lipid metabolism fractions in the animal-differentiating core relative to all other groups.}
\label{fig:functional_categories}
\end{figure}

\section{Supplementary Figures}
\label{sec:supp_figs}

\noindent Supplementary figures include:
\begin{itemize}
  \item Figure S1 --- Functional category distribution (absolute counts)
  \item Figure S2 --- Chromosome plot of animal-differentiating SPDIs
  \item Figure S3 --- Gene heatmap (top mutated genes per group)
  \item Figure S4 --- KEGG enrichment bubble plots (animal vs.\ ancestral; animal vs.\ basal)
  \item Figure S5 --- Per-species UpSet plot of SPDI overlaps
  \item Figure S6 --- Per-species dN/dS proxy distribution
  \item Figure S7 --- Robustness analysis (Jaccard similarity, containment matrix)
  \item Figure S8 --- PKS/PDIM pathway map (all species combined)
  \item Figure S9 --- LoF detail for convergently pseudogenised PKS genes
\end{itemize}

\end{document}
